\documentclass[10pt,twocolumn,a4paper]{article}

% --- Encoding / idioma (documento \'integramente en espa\~nol) ---
\usepackage[utf8]{inputenc}
\usepackage[T1]{fontenc}
\usepackage[english,spanish,es-noshorthands,es-nodecimaldot,es-nolists]{babel}

% --- Minimal margins for density ---
\usepackage[margin=0.7cm, top=1.2cm, bottom=1.2cm, columnsep=0.8cm]{geometry}

% --- Math ---
\usepackage{amsmath, amssymb, mathtools}
\numberwithin{equation}{subsection}

% --- Dense spacing ---
\usepackage{titlesec}
\usepackage{enumitem}

\setlength{\parindent}{0pt}
\setlength{\parskip}{3pt}

% --- Compact sections ---
\titleformat{\section}{\normalsize\bfseries}{\thesection.}{0.4em}{}
\titleformat{\subsection}{\small\bfseries}{\thesubsection.}{0.4em}{}
\titleformat{\subsubsection}{\small\itshape}{\thesubsubsection.}{0.4em}{}
\titlespacing*{\section}{0pt}{6pt}{2pt}
\titlespacing*{\subsection}{0pt}{4pt}{1pt}
\titlespacing*{\subsubsection}{0pt}{3pt}{1pt}

% --- Compact lists ---
\setlist{nosep, leftmargin=1.2em}

% --- Compact TOC ---
\usepackage{tocloft}
\setlength{\cftbeforesecskip}{1pt}
\setlength{\cftbeforesubsecskip}{0pt}
\renewcommand{\cftsecfont}{\small}
\renewcommand{\cftsubsecfont}{\footnotesize}
\renewcommand{\cftsecpagefont}{\small}
\renewcommand{\cftsubsecpagefont}{\footnotesize}

% --- Hyperref (loaded late to avoid conflicts) ---
\usepackage[hidelinks]{hyperref}

% --- float ---
\usepackage{float}

% --- tikz ---
\usepackage{tikz}
\usetikzlibrary{decorations.pathreplacing}

% --- Header ---
\usepackage{fancyhdr}
\pagestyle{fancy}
\fancyhf{}
\renewcommand{\headrulewidth}{0.4pt}
\fancyhead[L]{\footnotesize\leftmark}
\fancyhead[R]{\footnotesize\thepage}
\setlength{\headheight}{14pt}

% ============================================================
% CUSTOM MACROS
% ============================================================

% --- Shorthand: Greek letters ---
\newcommand{\al}{\alpha}
\newcommand{\be}{\beta}
\newcommand{\ga}{\gamma}
\newcommand{\Ga}{\Gamma}
\newcommand{\de}{\delta}
\newcommand{\De}{\Delta}
\newcommand{\ep}{\varepsilon}
\newcommand{\ze}{\zeta}
\newcommand{\et}{\eta}
\newcommand{\la}{\lambda}
\newcommand{\La}{\Lambda}
\newcommand{\si}{\sigma}
\newcommand{\Si}{\Sigma}
\newcommand{\ta}{\tau}
\newcommand{\om}{\omega}
\newcommand{\Om}{\Omega}
\newcommand{\ka}{\kappa}
\newcommand{\vph}{\varphi}
\newcommand{\ps}{\psi}
\newcommand{\xii}{\xi}

% --- Shorthand: Calculus & analysis ---
\newcommand{\dd}{\,\mathrm{d}}
\newcommand{\dv}[2]{\frac{\mathrm{d}#1}{\mathrm{d}#2}}
\newcommand{\ddv}[2]{\frac{\mathrm{d}^2#1}{\mathrm{d}#2^2}}
\newcommand{\dnv}[3]{\frac{\mathrm{d}^{#3}#1}{\mathrm{d}#2^{#3}}}
\newcommand{\pdv}[2]{\frac{\partial #1}{\partial #2}}
\newcommand{\pddv}[2]{\frac{\partial^2 #1}{\partial #2^2}}
\newcommand{\inte}[4]{\int_{#1}^{#2} #3 \dd #4}
\newcommand{\lito}{\mathrm{o}}              % o peque\~na de Landau

% --- Shorthand: Linear algebra ---
\newcommand{\vb}[1]{\mathbf{#1}}
\newcommand{\vh}[1]{\hat{\mathbf{#1}}}
\newcommand{\mat}[1]{\mathbf{#1}}
\newcommand{\tr}{\operatorname{tr}}
\newcommand{\diag}{\operatorname{diag}}
\newcommand{\rank}{\operatorname{rank}}
\newcommand{\spn}{\operatorname{span}}
\newcommand{\adj}{\operatorname{adj}}
\newcommand{\inner}[2]{\langle #1, #2 \rangle}
\newcommand{\braket}[2]{\langle #1 \,\vert\, #2 \rangle}   % producto interno de funciones
\newcommand{\norm}[1]{\lVert #1 \rVert}
\newcommand{\abs}[1]{\lvert #1 \rvert}
\newcommand{\tp}{\mathsf{T}}
\newcommand{\inv}{^{-1}}
\newcommand{\I}{\mat{I}}
\newcommand{\dg}{\dagger}

% --- Shorthand: Sets & logic ---
\newcommand{\R}{\mathbb{R}}
\newcommand{\Z}{\mathbb{Z}}
\newcommand{\N}{\mathbb{N}}
\newcommand{\Q}{\mathbb{Q}}
\newcommand{\C}{\mathbb{C}}
\newcommand{\F}{\mathbb{F}}
\newcommand{\Ee}{\mathbb{E}}
\newcommand{\such}{\mid}
\newcommand{\imp}{\Rightarrow}

% --- Shorthand: Operadores diferenciales ---
\DeclareMathOperator{\grad}{\mathbf{grad}}
\DeclareMathOperator{\dive}{\mathbf{div}}
\DeclareMathOperator{\rot}{\mathbf{rot}}
\newcommand{\Lap}{\nabla^{2}}
\newcommand{\Real}{\operatorname{Re}}
\newcommand{\Imag}{\operatorname{Im}}
\DeclareMathOperator*{\Sup}{\mathbf{sup}}
\newcommand{\Ext}[1]{\textstyle\bigwedge^{#1}\!L}   % espacio de #1-formas
\DeclareMathOperator{\Argu}{\mathbf{arg}}

% --- Sizing ---
\newcommand{\eval}[1]{\left. #1 \right|}
\newcommand{\qty}[1]{\left( #1 \right)}
\newcommand{\sqty}[1]{\left[ #1 \right]}
\newcommand{\set}[1]{\left\{ #1 \right\}}

% ============================================================
% ENTORNOS: tipo teorema (planos, sin cajas ni color)
% ============================================================
\newcounter{thm}[section]
\renewcommand{\thethm}{\thesection.\arabic{thm}}

\newenvironment{teorema}[1][]{%
  \refstepcounter{thm}\par\medskip\noindent
  \textbf{Teorema~\thethm\ifx\\#1\\\else\ (#1)\fi.}\itshape\quad
}{\par\medskip}

\newenvironment{lema}[1][]{%
  \refstepcounter{thm}\par\medskip\noindent
  \textbf{Lema~\thethm\ifx\\#1\\\else\ (#1)\fi.}\itshape\quad
}{\par\medskip}

\newenvironment{corolario}[1][]{%
  \refstepcounter{thm}\par\medskip\noindent
  \textbf{Corolario~\thethm\ifx\\#1\\\else\ (#1)\fi.}\itshape\quad
}{\par\medskip}

\newenvironment{proposicion}[1][]{%
  \refstepcounter{thm}\par\medskip\noindent
  \textbf{Proposici\'on~\thethm\ifx\\#1\\\else\ (#1)\fi.}\itshape\quad
}{\par\medskip}

\newenvironment{definicion}[1][]{%
  \refstepcounter{thm}\par\medskip\noindent
  \textbf{Definici\'on~\thethm\ifx\\#1\\\else\ (#1)\fi.}\quad
}{\par\medskip}

\newenvironment{ejemplo}[1][]{%
  \refstepcounter{thm}\par\medskip\noindent
  \textbf{Ejemplo~\thethm\ifx\\#1\\\else\ (#1)\fi.}\quad
}{\par\medskip}

\newenvironment{observacion}[1][]{%
  \par\medskip\noindent
  \textit{Observaci\'on\ifx\\#1\\\else\ (#1)\fi.}\quad
}{\par\medskip}

\newenvironment{demostracion}[1][Demostraci\'on]{%
  \par\medskip\noindent
  \textit{#1.}\quad
}{\hfill$\square$\par\medskip}

% ============================================================
\begin{document}

\twocolumn[{%
  \begin{center}
    {\large\bfseries M\'etodos Matem\'aticos para la F\'isica}\\[2pt]
    {\small AMMF \quad|\quad Apuntes de clase}\\[2pt]
    {\small Angel Luna \quad | \quad \'Ultima compilaci\'on: \today}
  \end{center}
  \vspace{-4pt}
  \hrule
}]

% El \'indice va en el flujo a dos columnas: dentro de \twocolumn[...] no
% podr\'ia partirse entre p\'aginas.
\tableofcontents
\vspace{4pt}
\hrule
\vspace{6pt}

\section{Notaci\'on y convenciones}
\label{sec:notacion}

Estos apuntes cubren dos bloques: c\'alculo vectorial y geometr\'ia diferencial
(Secciones \ref{sec:espacios_vectoriales}--\ref{sec:formas_diferenciales}) y
variable compleja, series y an\'alisis de Fourier
(Secciones \ref{sec:variable_compleja}--\ref{sec:integracion_compleja}).

\begin{itemize}
  \item $V$ denota un espacio vectorial sobre el campo $\F$, con
        $\F=\R$ o $\F=\C$; $\Ee^n$ denota el espacio eucl\'ideo de
        dimensi\'on $n$.
  \item $\vb{a},\vb{b},\vb{c}$ son vectores y $\vh{e}_i$ los elementos de una
        base ortonormal. La transposici\'on se escribe $\vb{a}^\tp$ y la
        transpuesta conjugada $\vb{a}^\dg$.
  \item $\wedge$ es el \textbf{producto cu\~na} o exterior; en $\R^3$ coincide
        con el producto cruz. El producto punto se escribe $\vb{a}\cdot\vb{b}$
        y el producto interno abstracto $g(\vb{a},\vb{b})$ o
        $\braket{\vb{a}}{\vb{b}}$.
  \item $\otimes$ es el producto directo y $\oplus$ la suma directa.
  \item $\vb{r}(t)$ denota una curva y $s$ su longitud de arco; la prima
        $\vb{T}'$ siempre significa derivada respecto de $s$.
  \item $\vb{M}(\xi,\et)$ denota una superficie parametrizada y
        $\vb{M}(\xi,\et,\zeta)$ un mapeo de volumen.
  \item En la parte compleja, $z=x+iy$, $\bar z$ es el conjugado,
        $\abs{z}=(z\bar z)^{1/2}$ y $\vph=\Argu(z)$ el argumento principal.
        $\Om$ y $D$ denotan dominios abiertos de $\C$ y $\ga$ una curva
        parametrizada contenida en ellos.
  \item $\lito(\vb{u})$ es la $\lito$ peque\~na de Landau, definida por
        $\lim_{\vb{u}\rightarrow\vb{0}}\norm{\lito(\vb{u})}/\norm{\vb{u}}=0$.
\end{itemize}

\subsection{Notaci\'on de Einstein}
\label{subsec:einstein}

Un \'indice repetido dentro de un mismo t\'ermino indica suma sobre todo su
rango:
\begin{equation}
  \label{eq:einstein}
  \sum_{i=1}^{n}a_ib_i=a_1b_1+a_2b_2+\dots+a_nb_n=a_ib_i .
\end{equation}
As\'i, la serie de Taylor $f(x)=\sum_{n=0}^{\infty}a_n(x-x_0)^n$ se abrevia
$f(x)=a_n(x-x_0)^n$.

\subsection{Delta de Kronecker}
\label{subsec:kronecker}

Sea $\set{\vh{e}_i}_{i=1}^{n}$ una base ortonormal de $\Ee^n$, es decir,
$\vh{e}_k$ tiene un $1$ en la entrada $k$-\'esima y ceros en las dem\'as.

\begin{definicion}[Delta de Kronecker]
  \label{def:kronecker}
  \begin{equation}
    \label{eq:kronecker}
    \de_{ij}=g(\vh{e}_i,\vh{e}_j)=\pdv{x_j}{x_i}=
    \begin{cases}
      0, & i\neq j,\\
      1, & i=j .
    \end{cases}
  \end{equation}
\end{definicion}

Su utilidad es \emph{contraer} sumas. Si $\vb{a}=\sum_{i=1}^{n}a_i\vh{e}_i$,
entonces
\begin{equation}
  \label{eq:contraccion_kronecker}
  \begin{aligned}
    g(\vb{a},\vh{e}_j)
      &=\qty{\sum_{i=1}^{n}a_i\vh{e}_i}\cdot\vh{e}_j
       =\sum_{i=1}^{n}a_i\,\vh{e}_i\cdot\vh{e}_j\\
      &=\sum_{i=1}^{n}a_i\de_{ij}=a_j ,
  \end{aligned}
\end{equation}
ya que todos los sumandos con $i\neq j$ se anulan.

\subsection{S\'imbolos de Levi-Civita}
\label{subsec:levi_civita}

Tres vectores linealmente independientes generan una caja de volumen
\begin{equation}
  \label{eq:volumen_caja}
  \mathrm{Vol}=(\text{\'area de la base})(\text{altura})
    =(\vb{a}\wedge\vb{b})\cdot\vb{c},
\end{equation}
cuyo signo define la \textbf{orientaci\'on} del sistema: es positiva si
$(\vb{a}\wedge\vb{b})\cdot\vb{c}>0$ y negativa en caso contrario.

\begin{definicion}[S\'imbolo de Levi-Civita]
  \label{def:levi_civita}
  Sea $\set{\vh{e}_i}_{i=1}^{3}$ la base can\'onica de $\Ee^3$. Entonces
  \begin{equation}
    \label{eq:levi_civita}
    \ep_{ijk}=(\vh{e}_i\wedge\vh{e}_j)\cdot\vh{e}_k=
    \begin{cases}
      0, & \text{si se repite alg\'un \'indice},\\
      -1, & \text{si la permutaci\'on es impar},\\
      1, & \text{si la permutaci\'on es par}.
    \end{cases}
  \end{equation}
\end{definicion}

Con \'el, la componente $i$-\'esima del producto cu\~na en $\R^3$ se escribe en
una sola l\'inea:
\begin{equation}
  \label{eq:cuna_levicivita}
  (\vb{a}\wedge\vb{b})_i=\ep_{ijk}a_jb_k .
\end{equation}

\begin{ejemplo}
  \label{ej:cuna_componentes}
  Para $i=1$ los \'unicos t\'erminos no nulos son $\ep_{123}$ y $\ep_{132}$:
  \begin{equation}
    \label{eq:cuna_primera_componente}
    \begin{aligned}
      (\vb{a}\wedge\vb{b})_1
        &=\ep_{123}a_2b_3+\ep_{132}a_3b_2\\
        &=a_2b_3-a_3b_2 ,
    \end{aligned}
  \end{equation}
  y an\'alogamente para las otras dos, de modo que
  \begin{equation}
    \label{eq:cuna_3d}
    \vb{a}\wedge\vb{b}=
    \begin{pmatrix}
      a_2b_3-a_3b_2\\
      a_3b_1-a_1b_3\\
      a_1b_2-a_2b_1
    \end{pmatrix}.
  \end{equation}
\end{ejemplo}

El triple producto y el rotacional admiten la misma abreviatura:
\begin{equation}
  \label{eq:triple_producto}
  (\vb{a}\wedge\vb{b})\cdot\vb{c}=\ep_{ijk}c_ia_jb_k,
  \qquad
  (\rot\vb{f})_i=\ep_{ijk}\partial_jf_k .
\end{equation}

\section{Espacios vectoriales}
\label{sec:espacios_vectoriales}

\subsection{Definiciones b\'asicas}
\label{subsec:definiciones_ev}

\begin{definicion}[Espacio vectorial]
  \label{def:espacio_vectorial}
  Un \textbf{espacio vectorial lineal} sobre un campo $\F$ ($\R$ o $\C$) es un
  conjunto $V$ de objetos, denotados $\vb{a},\vb{b},\dots$ y llamados
  \emph{vectores}, tal que a cada par $\vb{a},\vb{b}\in V$ le corresponde un
  vector $\vb{a}+\vb{b}\in V$ que cumple las propiedades usuales de la suma
  (conmutatividad, asociatividad, neutro e inverso) y a cada
  $\al\in\F$, $\vb{a}\in V$ le corresponde $\al\vb{a}\in V$ compatible con
  ellas.
\end{definicion}

\begin{definicion}[Base]
  \label{def:base}
  Una \textbf{base} de $V$ es un conjunto $B$ de elementos linealmente
  independientes que generan $V$. Una base finita genera un espacio de
  dimensi\'on finita y una infinita un espacio de dimensi\'on infinita.
\end{definicion}

Si $B=\set{\vb{a}_i}_{i=1}^{n}$ es un conjunto de elementos linealmente
independientes, hablamos de un espacio de dimensi\'on $n$. La independencia
lineal significa que
\begin{equation}
  \label{eq:independencia_lineal}
  \sum_{i=1}^{n}\al_i\vh{e}_i=\vb{0}
  \qquad\iff\qquad
  \al_i=0\ \ \forall i .
\end{equation}

\begin{teorema}
  \label{teo:dimension}
  Todas las bases de un espacio de dimensi\'on finita tienen el mismo n\'umero
  de vectores linealmente independientes.
\end{teorema}

\begin{ejemplo}[Espacios finitos]
  \label{ej:espacios_finitos}
  $\R^2$, $\R^3$ y en general $\R^n$ son espacios vectoriales de dimensi\'on
  $2$, $3$ y $n$ respectivamente, con $n\in\Z^{+}\cup\set{0}$.
\end{ejemplo}

\begin{ejemplo}[Base infinita]
  \label{ej:base_infinita}
  Sea $V$ un espacio de funciones y $f\in V$ continua por pedazos y acotada. En
  un punto interior del dominio donde $f$ est\'a bien definida podemos
  aproximarla en una vecindad abierta de $x_0$ mediante su serie de Taylor,
  \begin{equation}
    \label{eq:taylor_1d}
    f(x)=\sum_{n=0}^{\infty}\frac{1}{n!}
      \eval{\dnv{f}{x}{n}}_{x_0}(x-x_0)^n ,
  \end{equation}
  es decir, por una aproximaci\'on polinomial. Para este curso una base de ese
  espacio de funciones es
  \begin{equation}
    \label{eq:base_polinomial}
    B=\set{x^n}_{n=0}^{\infty}.
  \end{equation}
\end{ejemplo}

\subsection{Representaci\'on geom\'etrica}
\label{subsec:representacion_geometrica}

Sean $\vb{a},\vb{b}\in V$; entonces $\vb{c}=\vb{a}+\vb{b}\in V$. Dada una base
$B=\set{\vh{e}_i}_{i=1}^{n}$ se tiene
\begin{equation}
  \label{eq:expansion_base}
  \vb{a}=\sum_{i=1}^{n}a_i\vh{e}_i,
  \qquad
  \vb{b}=\sum_{i=1}^{n}b_i\vh{e}_i,
\end{equation}
de modo que la suma y el producto por escalar act\'uan componente a componente,
\begin{equation}
  \label{eq:suma_componentes}
  \vb{a}+\vb{b}=\sum_{i=1}^{n}(a_i+b_i)\vh{e}_i,
  \qquad
  \al\vb{a}=\sum_{i=1}^{n}(\al a_i)\vh{e}_i ,
\end{equation}
y el vector se representa por su columna de coordenadas en esa base,
\begin{equation}
  \label{eq:columna_coordenadas}
  \vb{a}=
  \begin{bmatrix}
    a_1\\a_2\\\vdots\\a_n
  \end{bmatrix}_{e}.
\end{equation}

Para espacios de dimensi\'on infinita la idea es la misma. Si $f$ est\'a bien
definida en un dominio $D$ y admite expansi\'on de Taylor, la base
$B=\set{x^n}_{n=0}^{\infty}$ da \eqref{eq:taylor_1d}. Pero existen otras bases
$B=\set{g_n(x)}_{n=0}^{\infty}$ igualmente buenas, respecto de las cuales
\begin{equation}
  \label{eq:expansion_general}
  f(x)=\sum_{n=0}^{\infty}\hat{a}_n\,g_n(x).
\end{equation}
La Secci\'on~\ref{sec:fourier} explota esta libertad tomando exponenciales
complejas como base.

\subsection{Producto interno}
\label{subsec:producto_interno}

\begin{definicion}[Producto interno]
  \label{def:producto_interno}
  Un \textbf{producto interno} (o escalar) es una operaci\'on que asocia a dos
  elementos de $V$ un n\'umero del campo,
  \begin{equation}
    \label{eq:producto_interno}
    g:V\otimes V\longrightarrow\F,
    \qquad
    g(\vb{a},\vb{b})=\al\in\F .
  \end{equation}
\end{definicion}

\begin{ejemplo}[Representaci\'on en $\R^n$]
  \label{ej:pi_real}
  \begin{equation}
    \label{eq:pi_real}
    g(\vb{a},\vb{b})=\vb{a}\cdot\vb{b}=\vb{a}^\tp\vb{b}
      =a_1b_1+a_2b_2+\dots+a_nb_n=a_ib_i .
  \end{equation}
\end{ejemplo}

\begin{ejemplo}[Representaci\'on en $\C^n$]
  \label{ej:pi_complejo}
  \begin{equation}
    \label{eq:pi_complejo}
    g(\vb{a},\vb{b})=\vb{a}^\dg\vb{b}
      =\sqty{a_1,a_2,\dots,a_n}^{*}
      \begin{bmatrix}
        b_1\\b_2\\\vdots\\b_n
      \end{bmatrix},
  \end{equation}
  donde $*$ denota conjugaci\'on compleja.
\end{ejemplo}

\begin{observacion}[Par\'entesis complejo]
  Si $a\in\C$ lo representamos como $a=x+iy$ con $x,y\in\R$ e
  $i=\sqrt{-1}$, y definimos su \textbf{conjugado} $\bar{a}=a^{*}=x-iy$.
\end{observacion}

Las propiedades del producto interno complejo son:
\begin{enumerate}
  \item \textbf{Hermiticidad}: $g(\vb{a},\vb{b})=g(\vb{b},\vb{a})^{*}$.
  \item \textbf{Antilinealidad en la primera entrada}:
        \begin{equation}
          \label{eq:antilineal}
          g(\al\vb{a}+\be\vb{b},\vb{c})
            =\al^{*}g(\vb{a},\vb{c})+\be^{*}g(\vb{b},\vb{c}).
        \end{equation}
  \item \textbf{Linealidad en la segunda entrada}:
        \begin{equation}
          \label{eq:lineal}
          g(\vb{a},\de\vb{b}+\ga\vb{c})
            =\de\, g(\vb{a},\vb{b})+\ga\, g(\vb{a},\vb{c}).
        \end{equation}
\end{enumerate}

\begin{observacion}
  De la hermiticidad se sigue que $g(\vb{a},\vb{a})=g(\vb{a},\vb{a})^{*}$, es
  decir, $g(\vb{a},\vb{a})\in\R$ aun en el caso complejo. Adem\'as
  $g(\vb{a},\vb{a})\ge 0$, con igualdad si y s\'olo si $\vb{a}=\vb{0}$.
\end{observacion}

En el caso real la hermiticidad se reduce a simetr\'ia,
$g(\vb{a},\vb{b})=g(\vb{b},\vb{a})$, y la operaci\'on es lineal en ambas
entradas; por ejemplo
\begin{equation}
  \label{eq:bilineal_real}
  (\al\vb{a}+\be\vb{b})\cdot\vb{c}
    =\al\,\vb{a}\cdot\vb{c}+\be\,\vb{b}\cdot\vb{c}.
\end{equation}

\subsection{Norma}
\label{subsec:norma}

\begin{teorema}[Desigualdad de Schwarz]
  \label{teo:schwarz}
  Sean $\vb{a},\vb{b}\in V$ con un producto interno definido sobre $V$.
  Entonces
  \begin{equation}
    \label{eq:schwarz}
    \abs{g(\vb{a},\vb{b})}^2\le\norm{\vb{a}}^2\norm{\vb{b}}^2 .
  \end{equation}
\end{teorema}

\begin{definicion}[Norma]
  \label{def:norma}
  La \textbf{norma} de un elemento de $V$ con producto interno puede definirse a
  trav\'es del mismo producto interno, y cumple:
  \begin{itemize}
    \item $\norm{\vb{0}}=0$;
    \item $\norm{\vb{a}}\ge 0$, con $\norm{\vb{a}}=0\iff\vb{a}=\vb{0}$;
    \item $\norm{\al\vb{a}}=\abs{\al}\norm{\vb{a}}$ para todo $\al\in\C$;
    \item desigualdad del tri\'angulo,
          $\norm{\vb{a}+\vb{b}}\le\norm{\vb{a}}+\norm{\vb{b}}$.
  \end{itemize}
\end{definicion}

En particular se define la \textbf{norma $p$},
\begin{equation}
  \label{eq:norma_p}
  \norm{\vb{a}}_p=\qty{\sum_{i=1}^{n}\abs{a_i}^{p}}^{1/p},
\end{equation}
de la cual los casos relevantes son
\begin{equation}
  \label{eq:norma_2}
  \norm{\vb{a}}_2=\qty{\sum_{i=1}^{n}\abs{a_i}^{2}}^{1/2}
    =g(\vb{a},\vb{a})^{1/2}=(\vb{a}^\dg\vb{a})^{1/2}
\end{equation}
y
\begin{equation}
  \label{eq:norma_1}
  \norm{\vb{a}}_1=\sum_{i=1}^{n}\abs{a_i}.
\end{equation}

\subsection{Ortonormalizaci\'on de Gram--Schmidt}
\label{subsec:gram_schmidt}

Dado un conjunto linealmente independiente
$B=\set{\vb{v}_i}_{i=1}^{n}$ se puede construir a partir de \'el un conjunto
ortonormal. Se normaliza el primer vector,
\begin{equation}
  \label{eq:gs_primero}
  \vh{u}_1=\frac{\vb{v}_1}{\norm{\vb{v}_1}},
\end{equation}
se le resta a $\vb{v}_2$ su proyecci\'on sobre $\vh{u}_1$,
\begin{equation}
  \label{eq:gs_proyeccion}
  \tilde{\vb{v}}_2=\vb{v}_2-\qty{\vb{v}_2\cdot\vh{u}_1}\vh{u}_1,
\end{equation}
y se normaliza el resultado,
\begin{equation}
  \label{eq:gs_segundo}
  \vh{u}_2=\frac{\tilde{\vb{v}}_2}{\norm{\tilde{\vb{v}}_2}} .
\end{equation}
En general, habiendo construido $\vh{u}_1,\dots,\vh{u}_{k-1}$,
\begin{equation}
  \label{eq:gs_general}
  \tilde{\vb{v}}_k=\vb{v}_k-\sum_{j=1}^{k-1}
    \qty{\vb{v}_k\cdot\vh{u}_j}\vh{u}_j,
  \qquad
  \vh{u}_k=\frac{\tilde{\vb{v}}_k}{\norm{\tilde{\vb{v}}_k}} .
\end{equation}

\subsection{Producto cu\~na en $\R^3$}
\label{subsec:producto_cuna_3d}

\begin{definicion}[Producto cu\~na]
  \label{def:producto_cuna}
  Sean $\vb{a},\vb{b}\in V\subset\R^3$. Definimos la operaci\'on
  $\wedge:V\otimes V\rightarrow V$ tal que
  $\vb{a}\wedge\vb{b}=\vb{c}$, donde $\vb{c}$ es perpendicular a ambos
  vectores.
\end{definicion}

Sus propiedades son:
\begin{enumerate}
  \item $\vb{a}\wedge\vb{b}=\vb{0}$ si $\vb{b}\parallel\vb{a}$ o
        $\vb{b}=\vb{0}$;
  \item $\vb{a}\wedge\vb{b}=-\vb{b}\wedge\vb{a}$ (antisimetr\'ia);
  \item $\vb{a}\wedge(\al\vb{b}+\be\vb{c})
         =\al\,\vb{a}\wedge\vb{b}+\be\,\vb{a}\wedge\vb{c}$;
  \item $(\al\vb{a}+\be\vb{b})\wedge\vb{c}
         =\al\,\vb{a}\wedge\vb{c}+\be\,\vb{b}\wedge\vb{c}$;
  \item $(\al\vb{a})\wedge\vb{b}=\al\qty{\vb{a}\wedge\vb{b}}$.
\end{enumerate}
El sentido de $\vb{c}$ lo fija la regla de la mano derecha.

\subsubsection{Interpretaci\'on f\'isica}
\label{subsubsec:torque}

El momento de una fuerza es
\begin{equation}
  \label{eq:torque}
  \boldsymbol{\ta}=\vb{r}\wedge\vb{F},
\end{equation}
con magnitud
\begin{equation}
  \label{eq:magnitud_torque}
  \norm{\boldsymbol{\ta}}=\norm{\vb{r}\wedge\vb{F}}
    =r F_{\perp}=\norm{\vb{r}}\norm{\vb{F}}\sin\theta .
\end{equation}

En general
\begin{equation}
  \label{eq:norma_cuna}
  \norm{\vb{a}\wedge\vb{b}}=\norm{\vb{a}}\norm{\vb{b}}\sin\theta,
\end{equation}
que es dos veces el \'area del tri\'angulo generado por $\vb{a}$ y $\vb{b}$, o
equivalentemente el \'area del paralelogramo que determinan. Por eso
$\vb{c}=\vb{a}\wedge\vb{b}=\vb{A}$ se interpreta como el \textbf{vector de
\'area} de esa caja bidimensional: tiene magnitud igual al \'area y direcci\'on
normal a ella.

\begin{observacion}
  Toda superficie tiene entonces un campo de vectores de \'area, uno sobre cada
  parche, y ese campo dota al espacio de una orientaci\'on. Como consecuencia,
  para una curva simple cerrada se dice que se recorre en \textbf{sentido
  positivo} si el interior de la curva queda del lado izquierdo.
\end{observacion}

\section{Repaso de c\'alculo en una dimensi\'on}
\label{sec:repaso_calculo}

\subsection{Funciones}
\label{subsec:funciones}

\begin{definicion}[Rango]
  \label{def:rango}
  El \textbf{rango} de una funci\'on $f:X\rightarrow Y$, con
  $X,Y\subset\R$, es el conjunto de los elementos de $Y$ que son valores que
  toma $f$:
  \begin{equation}
    \label{eq:rango}
    \mathrm{Ran}(f)=\set{y\in Y\such y=f(x)\ \text{para alg\'un}\ x\in X}.
  \end{equation}
\end{definicion}

\begin{definicion}[Sobreyectiva]
  \label{def:sobreyectiva}
  $f:X\rightarrow Y$ es \textbf{sobre} o \textbf{sobreyectiva} si cada elemento
  de $Y$ es imagen de alg\'un $x\in X$, esto es, si $\mathrm{Ran}(f)=Y$.
\end{definicion}

\begin{definicion}[Inyectiva]
  \label{def:inyectiva}
  $f:X\rightarrow Y$ es \textbf{uno a uno} o \textbf{inyectiva} si dos
  elementos distintos de $X$ tienen im\'agenes distintas en $Y$; es decir, si
  para cualesquiera $x_1,x_2\in X$ con $x_1\neq x_2$ se cumple
  $f(x_1)\neq f(x_2)$.
\end{definicion}

\subsection{L\'imites}
\label{subsec:limites_1d}

\begin{definicion}[L\'imite]
  \label{def:limite_1d}
  La funci\'on $f$ tiende al \textbf{l\'imite} $\ell$ en $x_0$ si
  $\forall\ep>0$ existe $\de_\ep>0$ tal que para todo $x$ con
  $0<\abs{x-x_0}<\de_\ep$ se cumple $\abs{f(x)-\ell}<\ep$. Simb\'olicamente,
  \begin{equation}
    \label{eq:limite_1d}
    \lim_{x\rightarrow x_0}f(x)=\ell .
  \end{equation}
\end{definicion}

\begin{ejemplo}
  \label{ej:limite_sandwich}
  Calculemos
  \begin{equation}
    \label{eq:limite_ejemplo}
    \lim_{x\rightarrow\infty}\frac{x+\sin^3x}{5x+6}=\frac{1}{5}.
  \end{equation}
  Como $-1\le\sin^3x\le 1$, la funci\'on queda encajonada,
  \begin{equation}
    \label{eq:sandwich}
    \frac{x-1}{5x+6}\le f(x)\le\frac{x+1}{5x+6}.
  \end{equation}
  Tomando $\xii=1/x$,
  \begin{equation}
    \label{eq:cambio_variable}
    \frac{x-1}{5x+6}
      =\frac{1-\frac{1}{x}}{5+\frac{6}{x}}
      =\frac{1-\xii}{5+6\xii},
  \end{equation}
  lo que transforma el l\'imite en $\lim_{\xii\rightarrow 0}$. Entonces
  \begin{equation}
    \label{eq:estimacion}
    \abs{\frac{1-\xii}{5+6\xii}-\frac{1}{5}}
      =\abs{\frac{-11\xii}{25+30\xii}}<\abs{\xii}<\de_\ep
  \end{equation}
  para $\xii$ suficientemente peque\~no, y basta tomar $\de_\ep=\ep$.
\end{ejemplo}

\subsection{Funciones continuas}
\label{subsec:continuidad_1d}

Para una funci\'on cualquiera no necesariamente se cumple
$\lim_{x\rightarrow x_0}f(x)=f(x_0)$.

\begin{definicion}[Continuidad]
  \label{def:continuidad_1d}
  $f$ es \textbf{continua} en $x_0$ si
  \begin{equation}
    \label{eq:continuidad_1d}
    \lim_{x\rightarrow x_0}f(x)=f(x_0),
  \end{equation}
  y es continua en un dominio $D$ si esto vale para todo $x_0\in D$.
\end{definicion}

\begin{teorema}
  \label{teo:algebra_continuas}
  Si $f$ y $g$ son continuas en $x_0$, entonces:
  \begin{enumerate}
    \item $f+g$ es continua en $x_0$;
    \item $fg$ es continua en $x_0$;
    \item si $g(x_0)\neq 0$, entonces $1/g$ es continua en $x_0$.
  \end{enumerate}
\end{teorema}

\begin{teorema}[Valor intermedio]
  \label{teo:valor_intermedio}
  Si $f$ es continua en $[a,b]$ y $f(a)<0<f(b)$, entonces existe alg\'un
  $x_0\in[a,b]$ tal que $f(x_0)=0$.
\end{teorema}

\begin{teorema}[Acotamiento]
  \label{teo:acotamiento}
  Si $f$ es continua en $[a,b]$, entonces $f$ est\'a acotada superiormente en
  $[a,b]$; es decir, existe $N\in\R$ tal que $f(x)\le N$ para todo
  $x\in[a,b]$.
\end{teorema}

\begin{teorema}[Valor extremo]
  \label{teo:valor_extremo}
  Si $f$ es continua en $[a,b]$, entonces existe alg\'un $y\in[a,b]$ tal que
  $f(y)\ge f(x)$ para todo $x\in[a,b]$.
\end{teorema}

\subsection{Derivaci\'on}
\label{subsec:derivacion_1d}

La pendiente de la recta secante que une los puntos de abscisas $x$ y $x+h$ es
$\qty{f(x+h)-f(x)}/h$. La pendiente de la recta \emph{tangente} en
$\qty{x,f(x)}$ se obtiene en el l\'imite.

\begin{definicion}[Derivada]
  \label{def:derivada_1d}
  $f$ es \textbf{derivable} en $x_0$ si existe
  \begin{equation}
    \label{eq:derivada_1d}
    \dv{}{x}f(x)=\lim_{h\rightarrow 0}\frac{f(x+h)-f(x)}{h}.
  \end{equation}
\end{definicion}

\begin{teorema}
  \label{teo:derivable_continua}
  Si $f$ es derivable en $x_0$, entonces $f$ es continua en $x_0$.
\end{teorema}

\begin{demostracion}
  Basta escribir la diferencia como el cociente incremental multiplicado por
  $h$:
  \begin{equation}
    \label{eq:demo_continuidad}
    \begin{aligned}
      \lim_{h\rightarrow 0}\qty{f(x+h)-f(x)}
        &=\lim_{h\rightarrow 0}\frac{f(x+h)-f(x)}{h}\,h\\
        &=\qty{\dv{}{x}f(x)}\lim_{h\rightarrow 0}h=0 ,
    \end{aligned}
  \end{equation}
  puesto que el primer factor existe y es finito por hip\'otesis. Por lo tanto
  $\lim_{h\rightarrow 0}f(x+h)=f(x)$ en $x_0$.
\end{demostracion}

\subsubsection{Consecuencias}
\label{subsubsec:reglas_derivacion}

Sean $f$ y $g$ derivables en $x\in(a,b)$. Entonces
\begin{align}
  \label{eq:regla_suma}
  \dv{}{x}(f+g)&=\dv{}{x}f+\dv{}{x}g,\\
  \label{eq:regla_producto}
  \dv{}{x}(fg)&=\qty{\dv{}{x}f}g+\qty{\dv{}{x}g}f,
\end{align}
y si $g(x_0)\neq 0$,
\begin{equation}
  \label{eq:regla_reciproco}
  \dv{}{x}\qty{\frac{1}{g}}=-\qty{\dv{}{x}g}\frac{1}{g^2}.
\end{equation}

\begin{demostracion}[Demostraci\'on de \eqref{eq:regla_suma}]
  \begin{equation}
    \label{eq:demo_suma}
    \begin{aligned}
      \dv{}{x}(f+g)
        &=\lim_{h\rightarrow 0}\frac{f(x+h)+g(x+h)-f(x)-g(x)}{h}\\
        &=\lim_{h\rightarrow 0}\frac{f(x+h)-f(x)}{h}\\
        &\quad+\lim_{h\rightarrow 0}\frac{g(x+h)-g(x)}{h}\\
        &=\dv{}{x}f+\dv{}{x}g .
    \end{aligned}
  \end{equation}
\end{demostracion}

\begin{demostracion}[Demostraci\'on de \eqref{eq:regla_producto}]
  Se suma y se resta el t\'ermino cruzado $f(x+h)g(x)$:
  \begin{equation}
    \label{eq:demo_producto}
    \small
    \begin{aligned}
      \dv{}{x}(fg)
        &=\lim_{h\rightarrow 0}\frac{f(x+h)g(x+h)-f(x)g(x)}{h}\\
        &=\lim_{h\rightarrow 0}\frac{1}{h}
          \big(f(x+h)g(x+h)-f(x+h)g(x)\\
        &\qquad\qquad+f(x+h)g(x)-f(x)g(x)\big)\\
        &=\lim_{h\rightarrow 0}f(x+h)\frac{g(x+h)-g(x)}{h}\\
        &\quad+\lim_{h\rightarrow 0}\frac{f(x+h)-f(x)}{h}g(x)\\
        &=f(x)\dv{}{x}g(x)+\qty{\dv{}{x}f(x)}g(x).
    \end{aligned}
  \end{equation}
\end{demostracion}

\begin{demostracion}[Demostraci\'on de \eqref{eq:regla_reciproco}]
  \begin{equation}
    \label{eq:demo_reciproco}
    \begin{aligned}
      \dv{}{x}\qty{\frac{1}{g}}
        &=\lim_{h\rightarrow 0}\frac{1}{h}
          \qty{\frac{1}{g(x+h)}-\frac{1}{g(x)}}\\
        &=\lim_{h\rightarrow 0}\frac{g(x)-g(x+h)}{h\,g(x)g(x+h)}\\
        &=-\lim_{h\rightarrow 0}\frac{g(x+h)-g(x)}{h}
          \lim_{h\rightarrow 0}\frac{1}{g(x)g(x+h)}\\
        &=-\qty{\dv{}{x}g(x)}\frac{1}{g(x)^2}.
    \end{aligned}
  \end{equation}
\end{demostracion}

\subsubsection{Regla de la cadena}
\label{subsubsec:regla_cadena_1d}

Si $f$ y $g$ son derivables en $x_0$ y $k(x)=(f\circ g)(x)=f(g(x))$, entonces
\begin{equation}
  \label{eq:regla_cadena_1d}
  \dv{}{x}f(g(x))=\qty{\eval{\dv{f}{u}}_{u=g(x)}}\dv{}{x}g(x).
\end{equation}

\begin{ejemplo}[Derivada de $x^n$]
  \label{ej:derivada_potencia}
  Sea $f(x)=x^n$ con $n\in\N$. Por el binomio de Newton,
  \begin{equation}
    \label{eq:derivada_potencia}
    \begin{aligned}
      \dv{}{x}f(x)
        &=\lim_{h\rightarrow 0}\frac{(x+h)^n-x^n}{h}\\
        &=\lim_{h\rightarrow 0}\frac{1}{h}
          \qty{x^n+nx^{n-1}h+\dots+h^n-x^n}\\
        &=\lim_{h\rightarrow 0}\qty{nx^{n-1}+\mathcal{O}(h)}
         =nx^{n-1}.
    \end{aligned}
  \end{equation}
\end{ejemplo}

\subsection{Teoremas del valor medio}
\label{subsec:valor_medio}

\begin{definicion}[Extremo local]
  \label{def:extremo_local}
  Sea $f$ una funci\'on y $A$ un conjunto contenido en su dominio. Un punto
  $x\in A$ es un \textbf{m\'aximo} (o \textbf{m\'inimo}) \textbf{local} de $f$
  sobre $A$ si existe $\de>0$ tal que $x$ es un m\'aximo (o m\'inimo) de $f$
  sobre $A\cap\qty{x-\tfrac{\de}{2},\,x+\tfrac{\de}{2}}$.
\end{definicion}

\begin{teorema}[Rolle]
  \label{teo:rolle}
  Si $f$ es continua en $[a,b]$, derivable en $(a,b)$ y $f(a)=f(b)$, entonces
  existe $x_0\in(a,b)$ tal que
  \begin{equation}
    \label{eq:rolle}
    \eval{\dv{}{x}f(x)}_{x_0}=0 .
  \end{equation}
\end{teorema}

\begin{teorema}[Valor medio]
  \label{teo:valor_medio}
  Si $f$ es continua en $[a,b]$ y derivable en $(a,b)$, entonces existe
  $x_0\in(a,b)$ tal que
  \begin{equation}
    \label{eq:valor_medio}
    \eval{\dv{}{x}f(x)}_{x_0}=\frac{f(b)-f(a)}{b-a}.
  \end{equation}
\end{teorema}

\section{C\'alculo vectorial}
\label{sec:calculo_vectorial}

Consideramos ahora funciones
$\vb{f}:X\subset\R^n\rightarrow\R^m$, con $n,m\in\N$.

\begin{ejemplo}[Campo escalar]
  \label{ej:campo_escalar}
  $f:X\subset\R^2\rightarrow\R$ con $f(\vb{x})=z$.
\end{ejemplo}

\begin{ejemplo}[Transformaci\'on $\R^2\rightarrow\R^2$]
  \label{ej:transformacion}
  \begin{equation}
    \label{eq:transformacion_polar}
    \vb{T}(\rho,\vph)=
    \begin{pmatrix}
      \rho\cos\vph\\
      \rho\sin\vph
    \end{pmatrix}.
  \end{equation}
\end{ejemplo}

\begin{ejemplo}[Superficie $\R^2\rightarrow\R^3$]
  \label{ej:superficie_esfera}
  \begin{equation}
    \label{eq:esfera_param}
    \vb{M}(\theta,\vph)=
    \begin{pmatrix}
      R\sin\theta\cos\vph\\
      R\sin\theta\sin\vph\\
      R\cos\theta
    \end{pmatrix},
  \end{equation}
  cuyas componentes satisfacen $x^2+y^2+z^2=R^2$.
\end{ejemplo}

\begin{ejemplo}[Curva $\R\rightarrow\R^3$]
  \label{ej:curva_R3}
  \begin{equation}
    \label{eq:curva_R3}
    \vb{C}(t)=
    \begin{pmatrix}
      x(t)\\y(t)\\z(t)
    \end{pmatrix},
    \qquad
    \vb{C}:I\subset\R\rightarrow\R^3 .
  \end{equation}
\end{ejemplo}

\begin{definicion}[Curvas y superficies de nivel]
  \label{def:curvas_nivel}
  Sea $f:X\subset\R^2\rightarrow\R$ escalar. Sus \textbf{curvas de nivel} son
  \begin{equation}
    \label{eq:curva_nivel}
    N_k=\set{(x,y)\in\R^2\such f(x,y)=k},
    \qquad k\in\R .
  \end{equation}
  An\'alogamente, si $f:X\subset\R^3\rightarrow\R$, las \textbf{superficies de
  nivel} son $f(x,y,z)=k$ con $k$ constante.
\end{definicion}

\subsection{L\'imites}
\label{subsec:limites_vectoriales}

\begin{definicion}[L\'imite vectorial]
  \label{def:limite_vectorial}
  Sea $\vb{f}:X\subset\R^n\rightarrow\R^m$. Decimos que
  \begin{equation}
    \label{eq:limite_vectorial}
    \lim_{\vb{x}\rightarrow\vb{a}}\vb{f}(\vb{x})=\vb{L}
  \end{equation}
  si dado $\ep>0$ existe $\de_\ep>0$ tal que si $\vb{x}\in X$ y
  $0<\norm{\vb{x}-\vb{a}}<\de_\ep$, entonces
  $\norm{\vb{f}(\vb{x})-\vb{L}}<\ep$.
\end{definicion}

Intuitivamente: para $\norm{\vb{f}-\vb{L}}$ arbitrariamente peque\~no existe
$\norm{\vb{x}-\vb{a}}$ suficientemente peque\~no. Aqu\'i
\begin{equation}
  \label{eq:componentes_f}
  \vb{f}=
  \begin{pmatrix}
    f_1(x_1,\dots,x_n)\\
    f_2(x_1,\dots,x_n)\\
    \vdots\\
    f_m(x_1,\dots,x_n)
  \end{pmatrix}.
\end{equation}

\begin{ejemplo}[Un l\'imite que no existe]
  \label{ej:limite_direccional}
  Sea $f:\R^2\setminus\set{\vb{0}}\rightarrow\R$ con
  \begin{equation}
    \label{eq:f_direccional}
    f(x,y)=\frac{x^2-y^2}{x^2+y^2}.
  \end{equation}
  En coordenadas polares $x=\rho\cos\vph$, $y=\rho\sin\vph$,
  \begin{equation}
    \label{eq:limite_direccional}
    \begin{aligned}
      \lim_{\vb{x}\rightarrow\vb{0}}f(x,y)
        &=\lim_{\rho\rightarrow 0}
          \frac{\rho^2\cos^2\vph-\rho^2\sin^2\vph}{\rho^2}\\
        &=\cos^2\vph-\sin^2\vph ,
    \end{aligned}
  \end{equation}
  que depende del \'angulo, es decir, de la direcci\'on de acercamiento: vale
  $1$ sobre $y=0$, $-1$ sobre $x=0$ y $0$ sobre $y=\pm x$. Si supusi\'eramos
  $L=1$ y buscamos $\abs{f-1}<\ep$, al aproximarnos por la ruta $x=0$
  obtendr\'iamos $\abs{-1-1}=2$, que no es arbitrariamente peque\~no. El
  l\'imite no existe.
\end{ejemplo}

\begin{teorema}[Unicidad]
  \label{teo:unicidad_limite}
  Si un l\'imite existe, es \'unico: si
  $\lim_{\vb{x}\rightarrow\vb{a}}\vb{f}=\vb{L}$ y
  $\lim_{\vb{x}\rightarrow\vb{a}}\vb{f}=\vb{M}$, entonces $\vb{M}=\vb{L}$.
\end{teorema}

\begin{teorema}[Propiedades algebraicas]
  \label{teo:algebra_limites}
  Sean $\vb{F},\vb{G}:X\subset\R^n\rightarrow\R^m$ vectoriales,
  $f,g:X\subset\R^n\rightarrow\R$ escalares y $k\in\R$. Si
  $\lim_{\vb{x}\rightarrow\vb{a}}\vb{F}=\vb{L}$ y
  $\lim_{\vb{x}\rightarrow\vb{a}}\vb{G}=\vb{M}$, entonces
  \begin{align}
    \label{eq:limite_suma}
    \lim_{\vb{x}\rightarrow\vb{a}}(\vb{F}+\vb{G})&=\vb{L}+\vb{M},\\
    \label{eq:limite_escalar}
    \lim_{\vb{x}\rightarrow\vb{a}}k\vb{F}&=k\vb{L}.
  \end{align}
  Si adem\'as $\lim_{\vb{x}\rightarrow\vb{a}}f=L$ y
  $\lim_{\vb{x}\rightarrow\vb{a}}g=M$, entonces
  $\lim_{\vb{x}\rightarrow\vb{a}}(fg)=LM$, y si $M\neq 0$ y $g$ no se anula
  cerca de $\vb{a}$,
  \begin{equation}
    \label{eq:limite_cociente}
    \lim_{\vb{x}\rightarrow\vb{a}}\frac{f}{g}=\frac{L}{M}.
  \end{equation}
\end{teorema}

\begin{teorema}[L\'imites por componentes]
  \label{teo:limites_componentes}
  Sea $\vb{f}:X\subset\R^n\rightarrow\R^m$. Entonces
  $\lim_{\vb{x}\rightarrow\vb{a}}\vb{f}(\vb{x})=\vb{L}$ con
  $\vb{L}=(L_1,\dots,L_m)^\tp$ si y s\'olo si
  $\lim_{\vb{x}\rightarrow\vb{a}}f_i(\vb{x})=L_i$ para cada $i$.
\end{teorema}

\begin{ejemplo}
  \label{ej:limite_existe}
  Sea $f:\R^2\setminus\set{\vb{0}}\rightarrow\R$ con
  $f(x,y)=\qty{x^3+y^5}/\qty{x^2+y^2}$. En polares,
  \begin{equation}
    \label{eq:limite_polares}
    \begin{aligned}
      \lim_{\vb{x}\rightarrow\vb{0}}f(x,y)
        &=\lim_{\rho\rightarrow 0}
          \frac{\rho^3\cos^3\vph+\rho^5\sin^5\vph}{\rho^2}\\
        &=\lim_{\rho\rightarrow 0}
          \qty{\rho\cos^3\vph+\rho^3\sin^5\vph}=0
    \end{aligned}
  \end{equation}
  para todo $\vph$. Aqu\'i el l\'imite s\'i existe, y la cota es uniforme en la
  direcci\'on: dado $\ep>0$ basta $\rho<\ep/2$, pues
  $\abs{\rho\cos^3\vph+\rho^3\sin^5\vph}\le\rho+\rho^3$.
\end{ejemplo}

\subsection{Continuidad}
\label{subsec:continuidad_vectorial}

\begin{definicion}[Continuidad]
  \label{def:continuidad_vectorial}
  Sea $\vb{f}:X\subset\R^n\rightarrow\R^m$ y $\vb{a}\in X$. Entonces $\vb{f}$
  es \textbf{continua} en $\vb{a}$ si $\vb{a}$ es un punto aislado o si
  \begin{equation}
    \label{eq:continuidad_vectorial}
    \lim_{\vb{x}\rightarrow\vb{a}}\vb{f}(\vb{x})=\vb{f}(\vb{a}).
  \end{equation}
  Si esto ocurre en todos los puntos de su dominio, $\vb{f}$ es continua en
  $X$.
\end{definicion}

\subsection{Derivadas parciales}
\label{subsec:derivadas_parciales}

\begin{definicion}[Derivada parcial]
  \label{def:derivada_parcial}
  La derivada parcial de $\vb{f}$ respecto de $x_i$ es la derivada de la
  funci\'on parcial correspondiente:
  \begin{equation}
    \label{eq:derivada_parcial}
    \small
    \pdv{\vb{f}}{x_i}=\lim_{h\rightarrow 0}
      \frac{\vb{f}(x_1,\dots,x_i+h,\dots,x_n)
            -\vb{f}(x_1,\dots,x_i,\dots,x_n)}{h}.
  \end{equation}
\end{definicion}

\begin{ejemplo}
  \label{ej:parcial_escalar}
  Sea $f:\R^2\rightarrow\R$ con $f(x,y)=x^2y+\cos(x+y)$. Entonces
  \begin{equation}
    \label{eq:parcial_escalar}
    \small
    \begin{aligned}
      \pdv{f}{x}
        &=\lim_{h\rightarrow 0}\frac{f(x+h,y)-f(x,y)}{h}\\
        &=\lim_{h\rightarrow 0}\frac{(x+h)^2y-x^2y}{h}\\
        &\quad+\lim_{h\rightarrow 0}
          \frac{\cos(x+h+y)-\cos(x+y)}{h}\\
        &=2xy-\sin(x+y).
    \end{aligned}
  \end{equation}
\end{ejemplo}

\begin{ejemplo}
  \label{ej:parcial_vectorial}
  Sea $\vb{f}:\R^2\rightarrow\R^2$ con
  \begin{equation}
    \label{eq:f_vectorial}
    \vb{f}(x,y)=
    \begin{pmatrix}
      2x^2+y\\
      x^3\sin(xy)
    \end{pmatrix}.
  \end{equation}
  La derivada parcial act\'ua componente a componente,
  \begin{equation}
    \label{eq:parcial_vectorial}
    \pdv{\vb{f}}{y}=
    \begin{pmatrix}
      \pdv{}{y}\qty{2x^2+y}\\[2pt]
      \pdv{}{y}\qty{x^3\sin(xy)}
    \end{pmatrix}
    =
    \begin{pmatrix}
      1\\
      x^4\cos(xy)
    \end{pmatrix}.
  \end{equation}
\end{ejemplo}

\subsection{Diferenciabilidad}
\label{subsec:diferenciabilidad}

Sean $V\subset\R^n$ y $W\subset\R^m$ espacios normados y
$\vb{f}:V\rightarrow W$. La idea es aproximar $\vb{f}$ por una parte lineal
m\'as un residuo que se anula m\'as r\'apido que $\norm{\vb{r}}$:
\begin{equation}
  \label{eq:aproximacion_lineal}
  \vb{f}(\vb{r})=\vb{g}(\vb{r})+\lito(\vb{r}),
\end{equation}
donde la $\lito$ peque\~na se define por
\begin{equation}
  \label{eq:o_pequena}
  \lim_{\vb{r}\rightarrow\vb{0}}
    \frac{\lito(\vb{r})}{\norm{\vb{r}}}=\vb{0}\in\R^m .
\end{equation}

\subsubsection{Analog\'ia en una dimensi\'on}
\label{subsubsec:analogia_1d}

Para $g:\R\rightarrow\R$, decir que $g$ es diferenciable en $x$ equivale a
\begin{equation}
  \label{eq:diferenciable_1d}
  \lim_{h\rightarrow 0}\frac{1}{h}
    \qty{g(x+h)-g(x)-h\dv{}{x}g}=0,
\end{equation}
o, escrito de otro modo,
\begin{equation}
  \label{eq:diferenciable_1d_bis}
  g(x+h)=g(x)+h\dv{}{x}g(x)+\lito(h).
\end{equation}

\begin{definicion}[Diferenciabilidad]
  \label{def:diferenciabilidad}
  $\vb{f}:V\subset\R^n\rightarrow W\subset\R^m$ es \textbf{diferenciable} si
  existe una transformaci\'on lineal $D\vb{f}(\vb{r}):V\rightarrow W$ tal que
  \begin{equation}
    \label{eq:diferenciabilidad}
    \vb{f}(\vb{r}+\vb{u})
      =\vb{f}(\vb{r})+D\vb{f}(\vb{r})[\vb{u}]+\lito(\vb{u}).
  \end{equation}
\end{definicion}

\begin{teorema}
  \label{teo:unicidad_diferencial}
  Si $\vb{f}$ es diferenciable en $\vb{r}\in\R^n$, hay una \'unica
  transformaci\'on lineal $D\vb{f}(\vb{r}):\R^n\rightarrow\R^m$ tal que
  \begin{equation}
    \label{eq:unicidad_diferencial}
    \lim_{\xii\rightarrow 0}
      \abs{\frac{\vb{f}(\vb{r}+\xii\vb{u})-\vb{f}(\vb{r})}
                {\norm{\xii\vb{u}}}-D\vb{f}(\vb{r})[\vb{u}]}=0 .
  \end{equation}
\end{teorema}

De aqu\'i se obtiene la expresi\'on operativa: la \textbf{derivada
direccional}
\begin{equation}
  \label{eq:derivada_direccional}
  \begin{aligned}
    D\vb{f}(\vb{r})[\vb{u}]
      &=\lim_{\xii\rightarrow 0}\frac{1}{\xii\norm{\vb{u}}}
        \qty{\vb{f}(\vb{r}+\xii\vb{u})-\vb{f}(\vb{r})}\\
      &=\frac{1}{\norm{\vb{u}}}
        \eval{\dv{}{\xii}\vb{f}(\vb{r}+\xii\vb{u})}_{\xii=0}.
  \end{aligned}
\end{equation}

\begin{ejemplo}[Caso escalar]
  \label{ej:gradiente_escalar}
  Sea $f:\R^n\rightarrow\R$. Escribiendo
  $\et_i=x_i+\xii u_i$ y aplicando la regla de la cadena,
  \begin{equation}
    \label{eq:derivada_direccional_escalar}
    \small
    \begin{aligned}
      Df(\vb{r})[\vb{u}]
        &=\frac{1}{\norm{\vb{u}}}
          \eval{\dv{}{\xii}f(\et_1,\dots,\et_n)}_{\xii=0}\\
        &=\frac{1}{\norm{\vb{u}}}
          \eval{\sum_{i=1}^{n}\pdv{f}{\et_i}\dv{\et_i}{\xii}}_{\xii=0}.
    \end{aligned}
  \end{equation}
  Como $\lim_{\xii\rightarrow 0}\et_i=x_i$ y
  $\dv{\et_i}{\xii}=u_i$, resulta
  \begin{equation}
    \label{eq:gradiente_producto}
    Df(\vb{r})[\vb{u}]=\frac{1}{\norm{\vb{u}}}
      \sum_{i=1}^{n}\pdv{f}{x_i}u_i
      =\nabla f\cdot\vh{u}.
  \end{equation}
\end{ejemplo}

\subsection{El gradiente}
\label{subsec:gradiente}

Generalizando el ejemplo anterior a $\vb{f}:\R^n\rightarrow\R^m$, la regla de
la cadena da
\begin{equation}
  \label{eq:jacobiano}
  \small
  D\vb{f}(\vb{r})[\vb{u}]=\frac{1}{\norm{\vb{u}}}
  \begin{pmatrix}
    \pdv{f_1}{r_1} & \dots & \pdv{f_1}{r_n}\\
    \pdv{f_2}{r_1} & \dots & \pdv{f_2}{r_n}\\
    \vdots & \ddots & \vdots\\
    \pdv{f_m}{r_1} & \dots & \pdv{f_m}{r_n}
  \end{pmatrix}
  \begin{pmatrix}
    u_1\\u_2\\\vdots\\u_n
  \end{pmatrix},
\end{equation}
es decir
\begin{equation}
  \label{eq:gradiente_matriz}
  D\vb{f}(\vb{r})[\vb{u}]=\grad\vb{f}\;\frac{\vb{u}}{\norm{\vb{u}}}
    =\grad\vb{f}\,\vh{u},
\end{equation}
donde la matriz de \eqref{eq:jacobiano} es el \textbf{gradiente} (o matriz
jacobiana) de $\vb{f}$.

\begin{ejemplo}
  \label{ej:gradiente_matriz}
  Sea $\vb{f}:\R^2\rightarrow\R^3$ con
  \begin{equation}
    \label{eq:f_ejemplo_gradiente}
    \vb{f}(x,y)=
    \begin{pmatrix}
      3x^2-y^3\\
      xy+2\\
      x-y+xy
    \end{pmatrix}
    \imp
    \grad\vb{f}=
    \begin{pmatrix}
      6x & -3y^2\\
      y & x\\
      1+y & -1+x
    \end{pmatrix}.
  \end{equation}
\end{ejemplo}

\begin{teorema}[Direcci\'on de m\'aximo cambio]
  \label{teo:maximo_cambio}
  Sea $f:\R^n\rightarrow\R$ con $\nabla f\neq\vb{0}$. Entonces
  $\nabla f(\vb{r})$ apunta en la direcci\'on en que $f$ tiene mayor cambio.
\end{teorema}

\begin{ejemplo}
  \label{ej:gradiente_silla}
  Sea $f:\R^2\rightarrow\R$ con $f=x^2-y^2$. Su gradiente es
  \begin{equation}
    \label{eq:gradiente_silla}
    \nabla f=
    \begin{pmatrix}
      \pdv{}{x}\qty{x^2-y^2}\\[2pt]
      \pdv{}{y}\qty{x^2-y^2}
    \end{pmatrix}
    =2
    \begin{pmatrix}
      x\\-y
    \end{pmatrix},
  \end{equation}
  y sobre la recta $y=x$ vale $\nabla f=2x(1,-1)^\tp$.
\end{ejemplo}

\subsection{Plano tangente}
\label{subsec:plano_tangente}

\begin{teorema}[El gradiente es normal a la superficie de nivel]
  \label{teo:gradiente_normal}
  Sea $f:\R^3\rightarrow\R$ un mapeo de clase $C^1$ y $S$ la superficie de
  nivel definida por $f(x,y,z)=k$. Si $\vb{v}$ es el vector tangente en
  $\xii=0$ a una curva $\vb{c}(\xii)$ contenida en $S$ con
  $\vb{c}(0)=(x_0,y_0,z_0)$, entonces
  $\eval{\nabla f}_{(x_0,y_0,z_0)}\cdot\vb{v}=0$.
\end{teorema}

\begin{demostracion}
  La curva vive sobre $f(x,y,z)=k$, luego
  $f(x(\xii),y(\xii),z(\xii))=k$ es constante y su derivada se anula. Por la
  regla de la cadena,
  \begin{equation}
    \label{eq:demo_gradiente_normal}
    \begin{aligned}
      \dv{}{\xii}f
        &=\pdv{f}{x}\dv{x}{\xii}+\pdv{f}{y}\dv{y}{\xii}
          +\pdv{f}{z}\dv{z}{\xii}\\
        &=\nabla f\cdot\dv{}{\xii}\vb{c}
         =\nabla f\cdot\vb{v}=0 .
    \end{aligned}
  \end{equation}
\end{demostracion}

\begin{ejemplo}
  \label{ej:esfera_nivel}
  Sea $f(x,y,z)=x^2+y^2+z^2$. La superficie de nivel $f=R^2>0$ es la esfera de
  radio $R$, y
  \begin{equation}
    \label{eq:gradiente_esfera}
    \nabla\qty{x^2+y^2+z^2}=2
    \begin{pmatrix}
      x\\y\\z
    \end{pmatrix},
  \end{equation}
  que en cada punto es radial, es decir, normal a la esfera.
\end{ejemplo}

Si dos curvas $\vb{c}_1$ y $\vb{c}_2$ contenidas en $S$ pasan por
$\vb{r}_0$ y sus tangentes
\begin{equation}
  \label{eq:tangentes_curvas}
  \vb{v}_1=\eval{\dv{}{\xii}\vb{c}_1}_{\vb{r}_0},
  \qquad
  \vb{v}_2=\eval{\dv{}{\xii}\vb{c}_2}_{\vb{r}_0}
\end{equation}
no son colineales, entonces son linealmente independientes y generan el
\textbf{plano tangente} a $S$ en $\vb{r}_0$. Como la ecuaci\'on de un plano de
normal $\vb{N}$ que pasa por $\vb{r}_0$ es
$\vb{N}\cdot(\vb{r}-\vb{r}_0)=0$, y por el
Teorema~\ref{teo:gradiente_normal} la normal es $\nabla f$, el plano tangente
es
\begin{equation}
  \label{eq:plano_tangente}
  \eval{\nabla f}_{\vb{r}_0}\cdot\qty{\vb{r}-\vb{r}_0}=0 .
\end{equation}

\subsection{Teorema de la funci\'on impl\'icita}
\label{subsec:funcion_implicita}

\begin{teorema}[Funci\'on impl\'icita]
  \label{teo:funcion_implicita}
  Sea $F:X\subset\R^n\rightarrow\R$ de clase $C^1$ y sea $\vb{a}$ un punto del
  conjunto de nivel
  $S=\set{\vb{x}\in\R^n\such F(\vb{x})=k}$. Si
  $\pdv{F}{x_i}(\vb{a})\neq 0$, entonces existen una vecindad $U$ de
  $(a_1,\dots,a_{i-1},a_{i+1},\dots,a_n)\in\R^{n-1}$, una vecindad $V$ de
  $a_i\in\R$ y una funci\'on $f:U\rightarrow V$ de clase $C^1$ tales que si
  $(x_1,\dots,x_{i-1},x_{i+1},\dots,x_n)\in U$ y $x_i\in V$ satisfacen
  $F(x_1,\dots,x_n)=k$, entonces
  $x_i=f(x_1,\dots,x_{i-1},x_{i+1},\dots,x_n)$.
\end{teorema}

\begin{ejemplo}
  \label{ej:elipsoide_implicita}
  Sea
  \begin{equation}
    \label{eq:elipsoide_nivel}
    f(x,y,z)=\frac{x^2}{4}+\frac{y^2}{36}+\frac{z^2}{9}=k .
  \end{equation}
  Tomando la superficie de nivel $k=1$ obtenemos un elipsoide. Como
  \begin{equation}
    \label{eq:parcial_elipsoide}
    \pdv{f}{z}=\frac{2z}{9},
  \end{equation}
  en el punto $\vb{a}=(0,0,-3)^\tp$, que pertenece al elipsoide, se tiene
  \begin{equation}
    \label{eq:parcial_elipsoide_a}
    \eval{\pdv{f}{z}}_{\vb{x}=\vb{a}}=-\frac{2}{3}\neq 0 .
  \end{equation}
  Por el teorema de la funci\'on impl\'icita, alrededor de ese punto podemos
  despejar
  \begin{equation}
    \label{eq:despeje_z}
    z=-3\sqrt{1-\frac{x^2}{4}-\frac{y^2}{36}} .
  \end{equation}
\end{ejemplo}

\begin{teorema}[Funci\'on impl\'icita generalizado]
  \label{teo:funcion_implicita_gen}
  Supongamos que $\vb{F}:A\subset\R^{n+m}\rightarrow\R^m$ es de clase $C^1$ y
  que $(\vb{a},\vb{b})\in A$ satisface $\vb{F}(\vb{a},\vb{b})=\vb{k}$. Si el
  determinante
  \begin{equation}
    \label{eq:determinante_implicita}
    \small
    \det
    \begin{pmatrix}
      \pdv{F_1}{y_1}(\vb{a},\vb{b}) & \dots & \pdv{F_1}{y_m}(\vb{a},\vb{b})\\
      \vdots & \ddots & \vdots\\
      \pdv{F_m}{y_1}(\vb{a},\vb{b}) & \dots & \pdv{F_m}{y_m}(\vb{a},\vb{b})
    \end{pmatrix}
    \neq 0,
  \end{equation}
  entonces hay una vecindad $U$ de $\vb{a}\in\R^n$ y una \'unica funci\'on
  $\vb{f}:U\rightarrow\R^m$ de clase $C^1$ tal que
  \begin{equation}
    \label{eq:conclusion_implicita}
    \vb{f}(\vb{a})=\vb{b},
    \qquad
    \vb{F}(\vb{x},\vb{f}(\vb{x}))=\vb{k}
    \quad\forall\vb{x}\in U .
  \end{equation}
\end{teorema}

\subsection{Teorema de la funci\'on inversa}
\label{subsec:funcion_inversa}

\begin{teorema}[Funci\'on inversa]
  \label{teo:funcion_inversa}
  Supongamos que $\vb{f}:\R^n\rightarrow\R^n$ es de clase $C^1$ sobre un
  abierto $A\subset\R^n$ y que en $\vb{a}\in A$ el jacobiano no se anula,
  \begin{equation}
    \label{eq:jacobiano_inversa}
    \abs{\frac{\partial(f_1,f_2,\dots,f_n)}
              {\partial(x_1,x_2,\dots,x_n)}}\neq 0 .
  \end{equation}
  Entonces existe un abierto $U\subset\R^n$ que contiene a $\vb{a}$ tal que
  $\vb{f}$ es uno a uno sobre $U$, el conjunto $V=\vb{f}(U)$ tambi\'en es
  abierto y hay una \'unica funci\'on inversa $\vb{g}:V\rightarrow U$ de
  $\vb{f}$, tambi\'en de clase $C^1$.
\end{teorema}

\section{Sistemas coordenados}
\label{sec:sistemas_coordenados}

\subsection{Dos dimensiones}
\label{subsec:coordenadas_2d}

\subsubsection{Cartesianas}
\label{subsubsec:cartesianas_2d}

\begin{equation}
  \label{eq:cartesianas_2d}
  \vb{r}=
  \begin{pmatrix}
    x\\y
  \end{pmatrix}
  =x\vh{x}+y\vh{y}=x\vh{e}_1+y\vh{e}_2 .
\end{equation}

\subsubsection{Polares}
\label{subsubsec:polares}

\begin{equation}
  \label{eq:polares}
  \vb{r}=
  \begin{pmatrix}
    \rho\\\vph
  \end{pmatrix}_{\vh{\rho},\vh{\vph}},
\end{equation}
con
\begin{equation}
  \label{eq:polares_directa}
  x=\rho\cos\vph,
  \qquad
  y=\rho\sin\vph,
\end{equation}
\begin{equation}
  \label{eq:polares_inversa}
  \rho=\qty{x^2+y^2}^{1/2},
  \qquad
  \vph=\arctan\frac{y}{x},
\end{equation}
y dominios $\rho\in[0,\infty)$, $\vph\in[0,2\pi)$.

\subsection{Tres dimensiones}
\label{subsec:coordenadas_3d}

\subsubsection{Cartesianas}
\label{subsubsec:cartesianas_3d}

\begin{equation}
  \label{eq:cartesianas_3d}
  \vb{r}=
  \begin{pmatrix}
    x\\y\\z
  \end{pmatrix}
  =x\vh{x}+y\vh{y}+z\vh{z}.
\end{equation}

\subsubsection{Cil\'indricas}
\label{subsubsec:cilindricas}

\begin{equation}
  \label{eq:cilindricas}
  \vb{r}=
  \begin{pmatrix}
    \rho\\\vph\\z
  \end{pmatrix}_{\vh{\rho},\vh{\vph},\vh{z}},
\end{equation}
con
\begin{equation}
  \label{eq:cilindricas_directa}
  \left\{
  \begin{aligned}
    x&=\rho\cos\vph,\\
    y&=\rho\sin\vph,\\
    z&=z,
  \end{aligned}
  \right.
  \qquad
  \left\{
  \begin{aligned}
    \rho&=\qty{x^2+y^2}^{1/2},\\
    \vph&=\arctan\qty{\frac{y}{x}},\\
    z&=z,
  \end{aligned}
  \right.
\end{equation}
y dominios
\begin{equation}
  \label{eq:cilindricas_dominios}
  \rho\in[0,\infty),
  \qquad
  \vph\in[0,2\pi),
  \qquad
  z\in\R .
\end{equation}

\subsubsection{Esf\'ericas}
\label{subsubsec:esfericas}

\begin{equation}
  \label{eq:esfericas}
  \vb{r}=
  \begin{pmatrix}
    r\\\theta\\\vph
  \end{pmatrix}_{\vh{r},\vh{\theta},\vh{\vph}},
\end{equation}
donde $\rho=r\sin\theta$ es el radio cil\'indrico, de modo que
\begin{equation}
  \label{eq:esfericas_directa}
  \left\{
  \begin{aligned}
    x&=\rho\cos\vph=r\sin\theta\cos\vph,\\
    y&=\rho\sin\vph=r\sin\theta\sin\vph,\\
    z&=r\cos\theta,
  \end{aligned}
  \right.
\end{equation}
con la transformaci\'on inversa
\begin{equation}
  \label{eq:esfericas_inversa}
  \left\{
  \begin{aligned}
    r&=\qty{x^2+y^2+z^2}^{1/2},\\
    \theta&=\arccos\frac{z}{r},\\
    \vph&=\arctan\frac{y}{x},
  \end{aligned}
  \right.
\end{equation}
y dominios $r\in[0,\infty)$, $\theta\in[0,\pi]$, $\vph\in[0,2\pi)$.

\section{Teor\'ia de curvas}
\label{sec:teoria_curvas}

\begin{definicion}[Curva]
  \label{def:curva}
  Una \textbf{curva} en $\R^n$ es una funci\'on
  $\vb{r}:I\subset\R\rightarrow\R^n$. Si $I=[a,b]$ con $a<b$, entonces
  $\vb{r}(a)$ y $\vb{r}(b)$ son los extremos de la curva.
\end{definicion}

\begin{ejemplo}[Recta]
  \label{ej:recta}
  \begin{equation}
    \label{eq:recta}
    \boldsymbol{\ell}(t)=\vb{r}_0+\vh{u}\,t,
    \qquad t\in\R .
  \end{equation}
\end{ejemplo}

\begin{ejemplo}[Elipse]
  \label{ej:elipse}
  Con $I=[0,2\pi]$,
  \begin{equation}
    \label{eq:elipse}
    \vb{r}(\vph)=
    \begin{pmatrix}
      a\cos\vph\\
      b\sin\vph
    \end{pmatrix}.
  \end{equation}
\end{ejemplo}

\begin{ejemplo}[H\'elice]
  \label{ej:helice}
  \begin{equation}
    \label{eq:helice}
    \vb{r}(t)=
    \begin{pmatrix}
      R\cos(\om t)\\
      R\sin(\om t)\\
      a t
    \end{pmatrix},
    \qquad t\in\R .
  \end{equation}
\end{ejemplo}

\begin{observacion}
  Una parametrizaci\'on arbitraria no siempre genera una curva
  \emph{decente}. Cuando la parametrizaci\'on s\'i produce una curva regular
  hablamos de un \textbf{mapeo}.
\end{observacion}

\subsection{Vector tangente}
\label{subsec:vector_tangente}

Definimos el vector tangente a $\vb{r}(t)$ como la velocidad instant\'anea
\begin{equation}
  \label{eq:velocidad}
  \vb{V}=\lim_{\De t\rightarrow 0}
    \frac{\vb{r}(t+\De t)-\vb{r}(t)}{\De t}=\dv{}{t}\vb{r}(t),
\end{equation}
y la rapidez instant\'anea como su norma,
\begin{equation}
  \label{eq:rapidez}
  V(t)=\norm{\vb{V}}=\norm{\dv{}{t}\vb{r}(t)} .
\end{equation}
Si $\vb{V}(t)$ es diferenciable, la aceleraci\'on es
\begin{equation}
  \label{eq:aceleracion}
  \vb{a}=\dv{}{t}\vb{V}=\ddv{}{t}\vb{r}(t).
\end{equation}

\subsection{Longitud de arco}
\label{subsec:longitud_arco}

Podemos aproximar la longitud de una curva por rectas que unen puntos
sucesivos. Partiendo $t\in[t_0,t_f]$ en $N$ trozos
$t_0,t_1,\dots,t_f$ con $\De t_i=t_{i+1}-t_i$,
\begin{equation}
  \label{eq:longitud_aproximada}
  S\approx\sum_{i=1}^{N}\norm{\vb{r}(t_i+\De t_i)-\vb{r}(t_i)} .
\end{equation}
Tomando $N\rightarrow\infty$, equivalentemente $\De t_i\rightarrow 0$,
\begin{equation}
  \label{eq:longitud_limite}
  S=\lim_{N\rightarrow\infty}\sum_{i=1}^{N}
    \norm{\frac{\vb{r}(t_i+\De t_i)-\vb{r}(t_i)}{\De t_i}}\De t_i,
\end{equation}
que en el l\'imite es la integral
\begin{equation}
  \label{eq:longitud_arco}
  S(t)=\int_{t_0}^{t}\norm{\dv{}{t'}\vb{r}(t')}\dd t' .
\end{equation}

Por el teorema fundamental del c\'alculo,
\begin{equation}
  \label{eq:derivada_arco}
  \dv{}{t}S(t)=\norm{\dv{}{t}\vb{r}(t)}>0
  \qquad\text{para}\quad t>t_0,
\end{equation}
de modo que $S$ es mon\'otonamente creciente. Vale entonces el teorema de la
funci\'on inversa y podemos escribir $t=t(s)$, es decir,
\textbf{parametrizar por longitud de arco}. Derivando la identidad
$t(s(t))=t$ por la regla de la cadena,
\begin{equation}
  \label{eq:inversa_arco}
  \dv{}{s}t=\frac{1}{\dv{}{t}s} .
\end{equation}

\begin{ejemplo}[Circunferencia]
  \label{ej:arco_circulo}
  Para $\vb{r}(\vph)=(R\cos\vph,\,R\sin\vph)^\tp$ con
  $\vph\in[0,\vph_1]$,
  \begin{equation}
    \label{eq:arco_circulo}
    \begin{aligned}
      S&=\int_{0}^{\vph_1}\norm{\dv{}{\vph}\vb{r}}\dd\vph\\
       &=\int_{0}^{\vph_1}\sqrt{\qty{\dv{x}{\vph}}^2
          +\qty{\dv{y}{\vph}}^2}\dd\vph=R\vph_1 .
    \end{aligned}
  \end{equation}
\end{ejemplo}

\begin{ejemplo}[H\'elice]
  \label{ej:arco_helice}
  Para \eqref{eq:helice},
  \begin{equation}
    \label{eq:arco_helice}
    \begin{aligned}
      \dv{}{t}\vb{r}(t)&=
      \begin{pmatrix}
        -R\om\sin(\om t)\\
        R\om\cos(\om t)\\
        a
      \end{pmatrix},\\
      S(t)&=\int_{0}^{t}\sqrt{R^2\om^2+a^2}\dd t'
           =\sqrt{R^2\om^2+a^2}\;t,
    \end{aligned}
  \end{equation}
  Escribamos $\La=\sqrt{R^2\om^2+a^2}$ para abreviar; entonces
  \begin{equation}
    \label{eq:t_de_s_helice}
    t=\frac{s}{\La},
  \end{equation}
  que describe la din\'amica de la curva. Claramente
  \begin{equation}
    \label{eq:derivadas_arco_helice}
    \dv{}{t}s=\La>0
    \imp
    \dv{}{s}t=\frac{1}{\La}>0 .
  \end{equation}
  Sustituyendo \eqref{eq:t_de_s_helice} en \eqref{eq:helice}, y con
  $\vph(s)=s\om/\La$, se obtiene la h\'elice parametrizada por longitud de
  arco,
  \begin{equation}
    \label{eq:helice_arco}
    \vb{r}(s)=
    \begin{pmatrix}
      R\cos\vph(s)\\
      R\sin\vph(s)\\
      a s/\La
    \end{pmatrix}.
  \end{equation}
\end{ejemplo}

\subsection{La triada de Fr\'enet--Serret}
\label{subsec:frenet_serret}

La parametrizaci\'on por longitud de arco permite construir una teor\'ia
robusta de curvas. Por la regla de la cadena,
\begin{equation}
  \label{eq:tangente_unitario}
  \dv{}{s}\vb{r}(s)=\dv{}{t}\vb{r}\,\dv{}{s}t
    =\frac{\vb{V}}{\norm{\vb{V}}}=\vb{T}(s),
\end{equation}
el \textbf{vector tangente unitario}, que cumple $\norm{\vb{T}}=1$.

De ah\'i se sigue una consecuencia inmediata: como
$\vb{T}\cdot\vb{T}=1$ es constante,
\begin{equation}
  \label{eq:T_perp_Tprima}
  \dv{}{s}\qty{\vb{T}\cdot\vb{T}}
    =\vb{T}'\cdot\vb{T}+\vb{T}\cdot\vb{T}'
    =2\,\vb{T}'\cdot\vb{T}=0
  \imp \vb{T}'\perp\vb{T},
\end{equation}
donde la prima denota derivada respecto de $s$.

\begin{definicion}[Normal y curvatura]
  \label{def:normal_curvatura}
  El \textbf{vector normal} $\vb{N}$ es el unitario perpendicular a $\vb{T}$
  definido por
  \begin{equation}
    \label{eq:normal}
    \vb{T}'=\ka(s)\vb{N},
    \qquad
    \ka(s)=\norm{\vb{T}'},
  \end{equation}
  y $\ka(s)$ es la \textbf{curvatura} de la curva.
\end{definicion}

\begin{definicion}[Binormal]
  \label{def:binormal}
  El \textbf{vector binormal} se define como
  \begin{equation}
    \label{eq:binormal}
    \vb{B}=\vb{T}\wedge\vb{N}.
  \end{equation}
  La terna $\set{\vb{T},\vb{N},\vb{B}}$ es el \textbf{marco de
  Fr\'enet--Serret}, local en cada punto de la curva.
\end{definicion}

Los tres vectores son unitarios,
$\norm{\vb{T}}=\norm{\vb{N}}=\norm{\vb{B}}=1$, y mutuamente ortogonales.

\subsection{Torsi\'on}
\label{subsec:torsion}

De $\vb{B}\cdot\vb{B}=1$ se obtiene, igual que en \eqref{eq:T_perp_Tprima},
\begin{equation}
  \label{eq:B_perp_Bprima}
  \dv{}{s}\qty{\vb{B}\cdot\vb{B}}=0
  \imp \vb{B}'\cdot\vb{B}=0
  \imp \vb{B}'\perp\vb{B}.
\end{equation}
Por otro lado, $\vb{B}\cdot\vb{T}=0$, de modo que
\begin{equation}
  \label{eq:B_perp_T}
  \begin{aligned}
    \dv{}{s}\qty{\vb{B}\cdot\vb{T}}&=0\\
    \imp \vb{B}'\cdot\vb{T}+\vb{B}\cdot\vb{T}'&=0 .
  \end{aligned}
\end{equation}
Pero $\vb{T}'=\ka(s)\vb{N}$ y $\vb{N}\perp\vb{B}$, luego el segundo t\'ermino
se anula y $\vb{B}'\cdot\vb{T}=0$, es decir $\vb{B}'\perp\vb{T}$. Siendo
$\vb{B}'$ perpendicular tanto a $\vb{B}$ como a $\vb{T}$, debe ser paralelo a
$\vb{N}$.

\begin{definicion}[Torsi\'on]
  \label{def:torsion}
  La \textbf{torsi\'on} $\ta(s)$ se define por
  \begin{equation}
    \label{eq:torsion}
    \vb{B}'=\ta(s)\vb{N},
    \qquad
    \abs{\ta(s)}=\norm{\vb{B}'} .
  \end{equation}
\end{definicion}

Falta la derivada de $\vb{N}$. Como $\vb{N}'\perp\vb{N}$, vive en el plano
generado por $\vb{T}$ y $\vb{B}$, es decir
$\vb{N}'=N_1\vb{T}+N_2\vb{B}$. Los coeficientes se obtienen derivando las
relaciones de ortogonalidad:
\begin{equation}
  \label{eq:coeficientes_Nprima}
  \begin{aligned}
    \vb{N}\cdot\vb{T}=0
      &\imp \vb{N}'\cdot\vb{T}+\vb{N}\cdot\vb{T}'=0\\
      &\imp \vb{N}'\cdot\vb{T}=-\ka(s),\\
    \vb{N}\cdot\vb{B}=0
      &\imp \vb{N}'\cdot\vb{B}+\vb{N}\cdot\vb{B}'=0\\
      &\imp \vb{N}'\cdot\vb{B}=-\ta(s).
  \end{aligned}
\end{equation}

\subsection{Recetario}
\label{subsec:recetario_frenet}

Reuniendo \eqref{eq:normal}, \eqref{eq:coeficientes_Nprima} y
\eqref{eq:torsion}, las \textbf{ecuaciones de Fr\'enet--Serret} son
\begin{equation}
  \label{eq:frenet_serret}
  \left\{
  \begin{aligned}
    \vb{T}'&=\ka(s)\vb{N},\\
    \vb{N}'&=-\ka(s)\vb{T}-\ta(s)\vb{B},\\
    \vb{B}'&=\ta(s)\vb{N}.
  \end{aligned}
  \right.
\end{equation}

Por el teorema de la funci\'on inversa se puede escribir la triada en
t\'erminos del par\'ametro original $t$, y entonces
\begin{align}
  \label{eq:curvatura_t}
  \ka(t)&=\frac{\norm{\dot{\vb{r}}(t)\wedge\ddot{\vb{r}}(t)}}
               {\norm{\dot{\vb{r}}(t)}^{3}},\\
  \label{eq:torsion_t}
  \ta(t)&=\frac{\dot{\vb{r}}(t)\cdot
                \qty{\ddot{\vb{r}}(t)\wedge\dddot{\vb{r}}(t)}}
               {\norm{\dot{\vb{r}}(t)\wedge\ddot{\vb{r}}(t)}^{2}},
\end{align}
donde
\begin{equation}
  \label{eq:puntos_derivadas}
  \dot{\vb{r}}=\dv{}{t}\vb{r}(t),
  \quad
  \ddot{\vb{r}}=\ddv{}{t}\vb{r}(t),
  \quad
  \dddot{\vb{r}}=\dnv{\vb{r}}{t}{3}(t).
\end{equation}

\begin{ejemplo}[Triada de la h\'elice]
  \label{ej:triada_helice}
  Con la notaci\'on de \eqref{eq:helice_arco}, derivando,
  \begin{equation}
    \label{eq:T_helice}
    \vb{T}=\dv{}{s}\vb{r}(s)=\frac{1}{\La}
    \begin{pmatrix}
      -R\om\sin\vph(s)\\
      R\om\cos\vph(s)\\
      a
    \end{pmatrix}.
  \end{equation}
  Derivando otra vez,
  \begin{equation}
    \label{eq:Tprima_helice}
    \vb{T}'=\underbrace{\frac{R\om^2}{R^2\om^2+a^2}}_{\ka(s)}
    \underbrace{
    \begin{pmatrix}
      -\cos\vph(s)\\
      -\sin\vph(s)\\
      0
    \end{pmatrix}}_{\vb{N}},
  \end{equation}
  de modo que la curvatura de la h\'elice es constante,
  $\ka=R\om^2/(R^2\om^2+a^2)$. El binormal resulta
  \begin{equation}
    \label{eq:B_helice}
    \vb{B}=\vb{T}\wedge\vb{N}=\frac{1}{\La}
    \begin{pmatrix}
      a\sin\vph(s)\\
      -a\cos\vph(s)\\
      R\om
    \end{pmatrix},
  \end{equation}
  y derivando,
  \begin{equation}
    \label{eq:Bprima_helice}
    \vb{B}'=\frac{a\om}{R^2\om^2+a^2}
    \begin{pmatrix}
      \cos\vph(s)\\
      \sin\vph(s)\\
      0
    \end{pmatrix}
    =-\frac{a\om}{R^2\om^2+a^2}\,\vb{N},
  \end{equation}
  es decir, la torsi\'on tambi\'en es constante,
  $\ta=-a\om/(R^2\om^2+a^2)$.
\end{ejemplo}

\section{Superficies parametrizadas}
\label{sec:superficies}

\begin{definicion}[Superficie parametrizada]
  \label{def:superficie}
  Una \textbf{superficie parametrizada} es un mapeo
  \begin{equation}
    \label{eq:superficie}
    \vb{M}:\Om\subset\R^2\longrightarrow\R^m,
    \qquad m=2,3,\dots
  \end{equation}
  En el caso $m=3$,
  \begin{equation}
    \label{eq:superficie_3d}
    \vb{M}(\xii,\et)=
    \begin{pmatrix}
      M_1(\xii,\et)\\
      M_2(\xii,\et)\\
      M_3(\xii,\et)
    \end{pmatrix}.
  \end{equation}
\end{definicion}

\begin{ejemplo}[Esfera]
  \label{ej:esfera_superficie}
  \begin{equation}
    \label{eq:esfera_superficie}
    \vb{M}(\theta,\vph)=
    \begin{pmatrix}
      R\sin\theta\cos\vph\\
      R\sin\theta\sin\vph\\
      R\cos\theta
    \end{pmatrix}.
  \end{equation}
\end{ejemplo}

\begin{ejemplo}[Cilindro]
  \label{ej:cilindro}
  \begin{equation}
    \label{eq:cilindro}
    \vb{M}(\vph,z)=
    \begin{pmatrix}
      R\cos\vph\\
      R\sin\vph\\
      z
    \end{pmatrix},
    \quad \vph\in[0,2\pi),\ z\in[-h,h].
  \end{equation}
\end{ejemplo}

\subsection{Diferencial de superficie}
\label{subsec:diferencial_superficie}

Sea $\vb{M}:\Om\subset\R^2\rightarrow\R^3$, de modo que
$\vb{M}=\vb{M}(\xii,\et)$. Si fijamos $\et=\et_0$ y tomamos el mapeo como
funci\'on de $\xii$, obtenemos
\begin{equation}
  \label{eq:curva_xi}
  \vb{M}_{\xii}=\vb{M}(\xii,\et_0),
\end{equation}
que es una curva sobre la superficie. An\'alogamente,
\begin{equation}
  \label{eq:curva_eta}
  \vb{M}_{\et}=\vb{M}(\xii_0,\et)
\end{equation}
es otra curva sobre $M$ en la direcci\'on $\et$. Los vectores tangentes a esas
curvas son
\begin{equation}
  \label{eq:tangentes_superficie}
  \vb{T}_{\xii}=\pdv{}{\xii}\vb{M}(\xii,\et_0),
  \qquad
  \vb{T}_{\et}=\pdv{}{\et}\vb{M}(\xii_0,\et),
\end{equation}
y la normal al plano tangente es su producto cu\~na,
\begin{equation}
  \label{eq:normal_superficie}
  \vb{N}=\eval{\pdv{}{\xii}\vb{M}\wedge\pdv{}{\et}\vb{M}}
    _{(\xii_0,\et_0)},
\end{equation}
de modo que el plano tangente en $\vb{M}_0$ es
$\vb{N}\cdot\qty{\vb{r}-\vb{M}_0}=0$.

La consecuencia importante es que las diferenciales de esas curvas, que viven
en el plano tangente, aproximan localmente a la superficie. En efecto, para
$\vb{M}(\xii(t),\et_0)$ y $\vb{M}(\xii_0,\et(t))$ la regla de la cadena da
\begin{equation}
  \label{eq:derivadas_curvas_superficie}
  \dv{}{t}\vb{M}_{\xii}=\pdv{}{\xii}\vb{M}\dv{\xii}{t},
  \qquad
  \dv{}{t}\vb{M}_{\et}=\pdv{}{\et}\vb{M}\dv{\et}{t},
\end{equation}
y por lo tanto
\begin{equation}
  \label{eq:diferenciales_curvas}
  \dd\vb{M}_{\xii}=\pdv{}{\xii}\vb{M}\dd\xii,
  \qquad
  \dd\vb{M}_{\et}=\pdv{}{\et}\vb{M}\dd\et .
\end{equation}

\begin{definicion}[Diferencial de superficie]
  \label{def:dS}
  \begin{equation}
    \label{eq:dS}
    \dd\vb{S}=\pdv{}{\xii}\vb{M}\wedge\pdv{}{\et}\vb{M}
      \dd\xii\dd\et ,
  \end{equation}
  y el \textbf{elemento de \'area} es su norma,
  \begin{equation}
    \label{eq:dS_norma}
    \dd S=\norm{\dd\vb{S}}
      =\norm{\pdv{}{\xii}\vb{M}\wedge\pdv{}{\et}\vb{M}}
       \abs{\dd\xii\dd\et} .
  \end{equation}
\end{definicion}

\subsection{Ejemplo: el elipsoide de revoluci\'on}
\label{subsec:elipsoide}

Sea $\vb{M}:\R^2\rightarrow\R^3$ con
\begin{equation}
  \label{eq:elipsoide_param}
  \vb{M}(\theta,\vph)=
  \begin{pmatrix}
    a\sin\theta\cos\vph\\
    a\sin\theta\sin\vph\\
    b\cos\theta
  \end{pmatrix}.
\end{equation}
Entonces
\begin{equation}
  \label{eq:dS_elipsoide}
  \small
  \begin{aligned}
    \dd\vb{S}
      &=\pdv{\vb{M}}{\theta}\wedge\pdv{\vb{M}}{\vph}
        \dd\theta\dd\vph\\
      &=
      \begin{pmatrix}
        a\cos\theta\cos\vph\\
        a\cos\theta\sin\vph\\
        -b\sin\theta
      \end{pmatrix}
      \wedge
      \begin{pmatrix}
        -a\sin\theta\sin\vph\\
        a\sin\theta\cos\vph\\
        0
      \end{pmatrix}
      \dd\theta\dd\vph\\
      &=
      \begin{pmatrix}
        ab\sin^2\theta\cos\vph\\
        ab\sin^2\theta\sin\vph\\
        a^2\sin\theta\cos\theta
      \end{pmatrix}
      \dd\theta\dd\vph\\
      &=a\sin\theta
      \begin{pmatrix}
        b\sin\theta\cos\vph\\
        b\sin\theta\sin\vph\\
        a\cos\theta
      \end{pmatrix}
      \dd\theta\dd\vph .
  \end{aligned}
\end{equation}

En particular, si $b=a$ recuperamos la esfera:
\begin{equation}
  \label{eq:dS_esfera}
  \dd\vb{S}=\vh{r}\,a^2\sin\theta\dd\theta\dd\vph,
  \qquad
  \vh{r}=
  \begin{pmatrix}
    \sin\theta\cos\vph\\
    \sin\theta\sin\vph\\
    \cos\theta
  \end{pmatrix}.
\end{equation}

\subsubsection{Integraci\'on del vector de \'area}
\label{subsubsec:integral_dS}

Integrando \eqref{eq:dS_elipsoide} sobre toda la superficie, las integrales en
$\vph$ de $\cos\vph$ y $\sin\vph$ se anulan y queda
\begin{equation}
  \label{eq:integral_dS}
  \begin{aligned}
    \iint\dd\vb{S}
      &=\int_{0}^{2\pi}\!\!\int_{0}^{\pi}
        \begin{pmatrix}
          b\sin\theta\cos\vph\\
          b\sin\theta\sin\vph\\
          a\cos\theta
        \end{pmatrix}
        a\sin\theta\dd\theta\dd\vph\\
      &=2\pi a^2\vh{z}\int_{0}^{\pi}\sin\theta\cos\theta\dd\theta
       =\vb{0},
  \end{aligned}
\end{equation}
como debe ser para toda superficie cerrada.

\subsubsection{Integral de \'area}
\label{subsubsec:integral_area}

\begin{equation}
  \label{eq:area_elipsoide}
  \begin{aligned}
    A&=\iint\norm{\dd\vb{S}}=\iint\dd S\\
     &=\int_{0}^{2\pi}\!\!\int_{0}^{\pi}
       a\sqrt{b^2\sin^2\theta+a^2\cos^2\theta}\,
       \sin\theta\dd\theta\dd\vph\\
     &=2\pi a\int_{0}^{\pi}
       \sqrt{b^2+(a^2-b^2)\cos^2\theta}\,\sin\theta\dd\theta .
  \end{aligned}
\end{equation}
En particular, si $a=b$,
\begin{equation}
  \label{eq:area_esfera}
  A=2\pi a^2\int_{0}^{\pi}\sin\theta\dd\theta=4\pi a^2,
\end{equation}
el \'area de la esfera.

\subsection{Diferencial de volumen}
\label{subsec:diferencial_volumen}

Supongamos ahora un mapeo $\vb{M}:\R^3\rightarrow\R^3$ con par\'ametros
$u,v,w$. Los vectores tangentes a las curvas coordenadas son
\begin{equation}
  \label{eq:tangentes_volumen}
  \vb{M}_u=\pdv{}{u}\vb{M},
  \qquad
  \vb{M}_v=\pdv{}{v}\vb{M},
  \qquad
  \vb{M}_w=\pdv{}{w}\vb{M} .
\end{equation}
Como el volumen de la caja que generan tres vectores es
$(\vb{a}\wedge\vb{b})\cdot\vb{c}$, el elemento de volumen es
\begin{equation}
  \label{eq:dV}
  \dd V=\qty{\vb{M}_u\wedge\vb{M}_v}\cdot\vb{M}_w
    \dd u\dd v\dd w .
\end{equation}

\begin{ejemplo}[Volumen del elipsoide]
  \label{ej:volumen_elipsoide}
  Tomemos el mapeo de la caja al elipsoide
  $\vb{M}=\vb{M}(\xii,\theta,\vph)$ con
  \begin{equation}
    \label{eq:mapeo_elipsoide}
    \vb{M}=
    \begin{pmatrix}
      a\xii\sin\theta\cos\vph\\
      b\xii\sin\theta\sin\vph\\
      c\xii\cos\theta
    \end{pmatrix},
    \qquad a<b<c .
  \end{equation}
  Un c\'alculo directo del triple producto da
  \begin{equation}
    \label{eq:dV_elipsoide}
    \dd V=\qty{\pdv{\vb{M}}{\xii}\wedge\pdv{\vb{M}}{\theta}}
      \cdot\pdv{\vb{M}}{\vph}\dd\xii\dd\theta\dd\vph
      =abc\,\xii^2\sin\theta\dd\xii\dd\theta\dd\vph ,
  \end{equation}
  de modo que
  \begin{equation}
    \label{eq:volumen_elipsoide}
    \begin{aligned}
      \mathrm{Vol}&=abc\int_{0}^{2\pi}\!\!\int_{0}^{\pi}\!\!\int_{0}^{1}
        \xii^2\sin\theta\dd\xii\dd\theta\dd\vph\\
        &=abc\,(2\pi)(2)\qty{\tfrac{1}{3}}=\frac{4\pi}{3}abc .
    \end{aligned}
  \end{equation}
  Si $a=b=c$ se recupera $\mathrm{Vol}=\tfrac{4}{3}\pi a^3$.
\end{ejemplo}

\begin{observacion}[Sistema positivo]
  Al construir $\dd\vb{S}$ o $\dd V$ hay que \emph{preservar el orden de las
  variables}: el signo del producto cu\~na y del triple producto depende de
  \'el, y con ello la orientaci\'on de la normal o el signo del volumen.
\end{observacion}

\section{Formas diferenciales y teoremas integrales}
\label{sec:formas_diferenciales}

\subsection{Teorema de Fubini}
\label{subsec:fubini}

\begin{teorema}[Fubini]
  \label{teo:fubini}
  Sea $f$ una funci\'on acotada en $D=[a,b]\otimes[c,d]$ y supongamos que el
  conjunto de discontinuidades $\Ga$ tiene \'area cero. Si cada l\'inea
  paralela a los ejes coordenados alcanza $\Ga$ en un n\'umero finito de
  puntos, entonces
  \begin{equation}
    \label{eq:fubini}
    \iint_{D}f\dd S
      =\int_{a}^{b}\!\!\int_{c}^{d}f(x,y)\dd y\dd x
      =\int_{c}^{d}\!\!\int_{a}^{b}f(x,y)\dd x\dd y .
  \end{equation}
\end{teorema}

\subsection{\'Algebra exterior}
\label{subsec:algebra_exterior}

\begin{definicion}[Cero-forma]
  \label{def:cero_forma}
  Una \textbf{cero-forma diferencial} es una funci\'on diferenciable
  $f:\R^n\rightarrow\R$. Escribimos $f\in\Ext{0}$, donde $\Ext{0}$ es el
  espacio de funciones sobre el espacio lineal $L$.
\end{definicion}

\begin{definicion}[Mapeo diferencial]
  \label{def:mapeo_diferencial}
  El \textbf{mapeo diferencial} $d$ act\'ua como
  $d:\Ext{n-1}\rightarrow\Ext{n}$. Sobre una cero-forma,
  \begin{equation}
    \label{eq:d_cero_forma}
    \mathbf{d}f=\pdv{f}{x_1}\dd x_1+\pdv{f}{x_2}\dd x_2
      +\dots+\pdv{f}{x_n}\dd x_n\in\Ext{1}.
  \end{equation}
\end{definicion}

\begin{definicion}[$n$-formas]
  \label{def:n_formas}
  Definimos el espacio de $n$-formas diferenciales como
  \begin{equation}
    \label{eq:espacio_n_formas}
    \Ext{n}=\underbrace{\Ext{1}\otimes\Ext{1}\otimes\dots\otimes\Ext{1}}
      _{n\ \text{factores}} ,
  \end{equation}
  cuyos elementos $\om\in\Ext{n}$ son
  \begin{equation}
    \label{eq:n_forma}
    \small
    \om=\sum_{i_1,\dots,i_n}
      \om_{i_1\dots i_n}(x_1,\dots,x_m)
      \dd x_{i_1}\wedge\dd x_{i_2}\wedge\dots\wedge\dd x_{i_n}.
  \end{equation}
\end{definicion}

\subsubsection{El producto cu\~na}
\label{subsubsec:producto_cuna_formas}

\begin{definicion}[Producto cu\~na de formas]
  \label{def:cuna_formas}
  \begin{equation}
    \label{eq:cuna_formas}
    \wedge:\Ext{p}\otimes\Ext{q}\longrightarrow\Ext{p+q},
  \end{equation}
  con las siguientes propiedades, si $\la\in\Ext{p}$ y $\mu\in\Ext{q}$:
  \begin{itemize}
    \item $\la\wedge\mu$ es distributivo;
    \item $\la\wedge(\mu\wedge\nu)=(\la\wedge\mu)\wedge\nu$, es decir, es
          asociativo;
    \item $\mu\wedge\la=(-1)^{pq}\la\wedge\mu$.
  \end{itemize}
\end{definicion}

De ah\'i se siguen las propiedades del \'algebra exterior:
\begin{align}
  \label{eq:exterior_dist_izq}
  (a_1\al_1+a_2\al_2)\wedge\be
    &=a_1\al_1\wedge\be+a_2\al_2\wedge\be,\\
  \label{eq:exterior_dist_der}
  \al\wedge(b_1\be_1+b_2\be_2)
    &=b_1\al\wedge\be_1+b_2\al\wedge\be_2,\\
  \label{eq:exterior_antisim}
  \al\wedge\be&=-\be\wedge\al .
\end{align}
En particular \eqref{eq:exterior_antisim} implica
$\dd x_i\wedge\dd x_i=0$, y por definici\'on $d(\dd x_i)=0$.

\begin{ejemplo}
  \label{ej:cuna_descompuesta}
  Si $\la=\la_1\wedge\dots\wedge\la_p$ y
  $\mu=\mu_1\wedge\dots\wedge\mu_q$, entonces
  \begin{equation}
    \label{eq:cuna_descompuesta}
    \la\wedge\mu
      =\la_1\wedge\dots\wedge\la_p\wedge\mu_1\wedge\dots\wedge\mu_q .
  \end{equation}
\end{ejemplo}

\begin{ejemplo}[Una $1$-forma por una $2$-forma]
  \label{ej:uno_por_dos}
  Sean $\la\in\Ext{1}$ y $\mu\in\Ext{2}$ en $\R^3$,
  \begin{equation}
    \label{eq:lambda_mu}
    \begin{aligned}
      \la&=\la_1\dd x+\la_2\dd y+\la_3\dd z,\\
      \mu&=\mu_1\dd y\wedge\dd z+\mu_2\dd z\wedge\dd x
          +\mu_3\dd x\wedge\dd y .
    \end{aligned}
  \end{equation}
  Todos los productos con un diferencial repetido se anulan y los tres que
  sobreviven se reordenan al mismo signo, de modo que
  \begin{equation}
    \label{eq:lambda_cuna_mu}
    \la\wedge\mu=\qty{\la_1\mu_1+\la_2\mu_2+\la_3\mu_3}
      \dd x\wedge\dd y\wedge\dd z .
  \end{equation}
\end{ejemplo}

\subsubsection{F\'ormula general}
\label{subsubsec:formula_general_cuna}

Para $k$ vectores $\vb{v}_1,\dots,\vb{v}_k\in V$, su producto exterior es
\begin{equation}
  \label{eq:cuna_general}
  \small
  \vb{v}_1\wedge\dots\wedge\vb{v}_k
    =\sum_{1\le i_1<\dots<i_k\le n}
    \begin{vmatrix}
      v_1^{i_1} & \dots & v_1^{i_k}\\
      \vdots & \ddots & \vdots\\
      v_k^{i_1} & \dots & v_k^{i_k}
    \end{vmatrix}
    \vh{e}_{i_1}\wedge\dots\wedge\vh{e}_{i_k}.
\end{equation}

\begin{ejemplo}
  \label{ej:cuna_general_R3}
  Calculemos el producto exterior de dos vectores de $\R^3$ en la base
  can\'onica:
  \begin{equation}
    \label{eq:cuna_general_R3}
    \small
    \begin{aligned}
      \vb{v}\wedge\vb{w}
        &=\sum_{1\le i<j\le 3}
        \begin{vmatrix}
          v_i & v_j\\
          w_i & w_j
        \end{vmatrix}
        \vh{e}_i\wedge\vh{e}_j\\
        &=(v_1w_2-v_2w_1)\,\vh{e}_1\wedge\vh{e}_2\\
        &\quad+(v_1w_3-v_3w_1)\,\vh{e}_1\wedge\vh{e}_3\\
        &\quad+(v_2w_3-v_3w_2)\,\vh{e}_2\wedge\vh{e}_3 ,
    \end{aligned}
  \end{equation}
  donde $\vh{e}_i\wedge\vh{e}_j\equiv\vh{e}_{ij}$. Identificando
  $\vh{e}_{23}\mapsto\vh{e}_1$, $\vh{e}_{31}\mapsto\vh{e}_2$ y
  $\vh{e}_{12}\mapsto\vh{e}_3$ se recupera el producto cruz
  \eqref{eq:cuna_3d}.
\end{ejemplo}

\subsection{La derivada exterior en $\R^3$}
\label{subsec:derivada_exterior}

Sea $\om\in\Ext{1}$ con $\om_i:\R^3\rightarrow\R$,
\begin{equation}
  \label{eq:uno_forma}
  \om=\om_1\dd x+\om_2\dd y+\om_3\dd z .
\end{equation}
Entonces, usando $d(\dd x_i)=0$,
\begin{equation}
  \label{eq:d_omega_paso1}
  d\om=d\om_1\wedge\dd x+d\om_2\wedge\dd y+d\om_3\wedge\dd z,
\end{equation}
y desarrollando cada $d\om_i$ con \eqref{eq:d_cero_forma},
\begin{equation}
  \label{eq:d_omega_paso2}
  \small
  \begin{aligned}
    d\om&=\qty{\pdv{\om_1}{x}\dd x+\pdv{\om_1}{y}\dd y
              +\pdv{\om_1}{z}\dd z}\wedge\dd x\\
        &\quad+\qty{\pdv{\om_2}{x}\dd x+\pdv{\om_2}{y}\dd y
              +\pdv{\om_2}{z}\dd z}\wedge\dd y\\
        &\quad+\qty{\pdv{\om_3}{x}\dd x+\pdv{\om_3}{y}\dd y
              +\pdv{\om_3}{z}\dd z}\wedge\dd z .
  \end{aligned}
\end{equation}
Los t\'erminos con diferencial repetido se anulan y los restantes se agrupan:
\begin{equation}
  \label{eq:d_omega_rot}
  \begin{aligned}
    d\om&=\qty{\pdv{\om_2}{x}-\pdv{\om_1}{y}}\dd x\wedge\dd y\\
        &\quad+\qty{\pdv{\om_1}{z}-\pdv{\om_3}{x}}\dd z\wedge\dd x\\
        &\quad+\qty{\pdv{\om_3}{y}-\pdv{\om_2}{z}}\dd y\wedge\dd z .
  \end{aligned}
\end{equation}

Ahora bien, el diferencial de superficie se descompone como
\begin{equation}
  \label{eq:dS_componentes}
  \left\{
  \begin{aligned}
    \dd\vb{S}_x&=\vh{x}\dd y\wedge\dd z,\\
    \dd\vb{S}_y&=\vh{y}\dd z\wedge\dd x,\\
    \dd\vb{S}_z&=\vh{z}\dd x\wedge\dd y,
  \end{aligned}
  \right.
\end{equation}
de modo que \eqref{eq:d_omega_rot} se lee como un producto punto:
\begin{equation}
  \label{eq:d_omega_producto}
  d\om=
  \underbrace{\qty{\nabla\wedge
    \begin{pmatrix}\om_1\\\om_2\\\om_3\end{pmatrix}}}_{\rot\boldsymbol{\om}}
  \cdot
  \underbrace{
  \begin{pmatrix}
    \dd y\wedge\dd z\\
    \dd z\wedge\dd x\\
    \dd x\wedge\dd y
  \end{pmatrix}}_{\dd\vb{S}} .
\end{equation}

\subsection{Teorema de Stokes}
\label{subsec:stokes}

\begin{teorema}[Stokes generalizado]
  \label{teo:stokes_general}
  \begin{equation}
    \label{eq:stokes_general}
    \int_{\partial S}\om=\int_{S}d\om .
  \end{equation}
\end{teorema}

Recordando que para $\om\in\Ext{1}$ se tiene
$\om=\boldsymbol{\om}\cdot\dd\vb{r}$, la combinaci\'on de
\eqref{eq:stokes_general} con \eqref{eq:d_omega_producto} da la forma
vectorial cl\'asica
\begin{equation}
  \label{eq:stokes_vectorial}
  \int_{C=\partial S}\vb{F}\cdot\dd\vb{r}
    =\iint_{S}\rot\vb{F}\cdot\dd\vb{S} .
\end{equation}

\begin{observacion}
  El teorema s\'olo es v\'alido en superficies \emph{abiertas}: una superficie
  cerrada no tiene frontera, de modo que el miembro izquierdo no existe.
\end{observacion}

\begin{observacion}[Invariancia de norma]
  Supongamos $\vb{G}=\rot\vb{F}$. Como el rotacional de un gradiente es cero,
  tambi\'en $\vb{G}=\rot\qty{\vb{F}+\nabla g}$ para cualquier campo escalar
  $g$ de clase $C^2$, y entonces
  \begin{equation}
    \label{eq:invariancia_norma}
    \iint_{S}\vb{G}\cdot\dd\vb{S}
      =\int_{\partial S}\qty{\vb{F}+\nabla g}\cdot\dd\vb{r},
  \end{equation}
  es decir, $\vb{F}$ est\'a determinado s\'olo salvo un gradiente.
\end{observacion}

\subsubsection{Ley de Amp\`ere}
\label{subsubsec:ampere}

En su forma integral,
\begin{equation}
  \label{eq:ampere_integral}
  \int_{C=\partial S}\vb{B}\cdot\dd\vb{r}=\mu_0 I .
\end{equation}
Escribiendo la corriente como el flujo de la densidad de corriente
$\vb{J}$ y aplicando \eqref{eq:stokes_vectorial} al miembro izquierdo,
\begin{equation}
  \label{eq:ampere_paso}
  \iint_{S}\rot\vb{B}\cdot\dd\vb{S}
    =\mu_0\iint_{S}\vb{J}\cdot\dd\vb{S},
\end{equation}
de donde
\begin{equation}
  \label{eq:ampere_diferencial}
  \iint_{S}\qty{\rot\vb{B}-\mu_0\vb{J}}\cdot\dd\vb{S}=0
  \quad\forall S .
\end{equation}
Como esto vale para \emph{toda} superficie, el integrando debe anularse:
\begin{equation}
  \label{eq:ampere_local}
  \rot\vb{B}=\mu_0\vb{J},
  \qquad\text{es decir}\qquad
  \nabla\wedge\vb{B}=\mu_0\vb{J} .
\end{equation}
Aqu\'i $\vb{B}=\rot\vb{A}$, con $\vb{A}$ el \textbf{potencial vectorial};
por la observaci\'on anterior, $\vb{A}$ y $\vb{A}+\nabla\la$ producen el mismo
campo, que es la libertad de norma del electromagnetismo.

\subsection{Teorema de Gauss}
\label{subsec:gauss}

Sea ahora $\om\in\Ext{2}$,
\begin{equation}
  \label{eq:dos_forma}
  \om=\om_1\dd y\wedge\dd z+\om_2\dd z\wedge\dd x
    +\om_3\dd x\wedge\dd y .
\end{equation}
Repitiendo el c\'alculo de \eqref{eq:d_omega_paso2}, ahora sobreviven
\'unicamente los t\'erminos que completan las tres coordenadas:
\begin{equation}
  \label{eq:d_omega_div}
  \begin{aligned}
    d\om&=\qty{\pdv{\om_1}{x}+\pdv{\om_2}{y}+\pdv{\om_3}{z}}
          \dd x\wedge\dd y\wedge\dd z\\
        &=\qty{\nabla\cdot\boldsymbol{\om}}\dd V .
  \end{aligned}
\end{equation}

\begin{teorema}[Divergencia]
  \label{teo:divergencia}
  \begin{equation}
    \label{eq:divergencia}
    \iint_{S=\partial V}\vb{F}\cdot\dd\vb{S}
      =\iiint_{V}\dive\vb{F}\dd V .
  \end{equation}
  Si la superficie no es cerrada, el teorema no aplica.
\end{teorema}

\begin{ejemplo}
  \label{ej:divergencia}
  Sea
  \begin{equation}
    \label{eq:F_divergencia}
    \vb{F}=
    \begin{pmatrix}
      3xy^2\\3yx^2\\z^3
    \end{pmatrix}
  \end{equation}
  y busquemos su flujo sobre la esfera unitaria. Al ser una superficie cerrada
  conviene el teorema de la divergencia:
  \begin{equation}
    \label{eq:div_F}
    \dive\vb{F}=\nabla\cdot\vb{F}=3y^2+3x^2+3z^2=3\qty{x^2+y^2+z^2},
  \end{equation}
  que en esf\'ericas es simplemente $\dive\vb{F}=3r^2$. Entonces
  \begin{equation}
    \label{eq:flujo_esfera}
    \begin{aligned}
      \iint_{S=\partial V}\vb{F}\cdot\dd\vb{S}
        &=\int_{0}^{2\pi}\!\!\int_{0}^{\pi}\!\!\int_{0}^{1}
          \qty{3r^2}r^2\sin\theta\dd r\dd\theta\dd\vph\\
        &=3(4\pi)\int_{0}^{1}r^4\dd r=\frac{12}{5}\pi .
    \end{aligned}
  \end{equation}
\end{ejemplo}

\section{Variable compleja}
\label{sec:variable_compleja}

\subsection{Representaci\'on polar}
\label{subsec:representacion_polar}

Sea $z=x+iy=(x,y)$, referido a la base $\set{1,i}$ del plano complejo.
Definimos el valor absoluto de $z$ como
\begin{equation}
  \label{eq:modulo}
  \abs{z}=\qty{z\bar{z}}^{1/2}=\abs{\bar{z}}=\qty{x^2+y^2}^{1/2},
\end{equation}
lo que permite la \textbf{representaci\'on polar}
\begin{equation}
  \label{eq:polar_compleja}
  z=\abs{z}\qty{\cos\vph+i\sin\vph},
  \qquad
  \Argu(z)=\vph=\arctan\qty{\frac{y}{x}} .
\end{equation}

\begin{observacion}
  Bajo esta representaci\'on el argumento no es \'unico:
  \begin{equation}
    \label{eq:argumento_multivaluado}
    z=\abs{z}\qty{\cos(\vph+2\pi n)+i\sin(\vph+2\pi n)},
    \quad n\in\Z .
  \end{equation}
  Se llama \textbf{argumento principal} al que cumple $\vph\in[0,2\pi)$.
\end{observacion}

\begin{ejemplo}
  \label{ej:polar}
  Para $z=2+i$ se tiene
  $z=\sqrt{5}\qty{\cos\al+i\sin\al}$ con
  $\al=\arctan\qty{\tfrac{1}{2}}$.
\end{ejemplo}

\subsection{Potencias}
\label{subsec:potencias}

Elevando al cuadrado \eqref{eq:polar_compleja} y usando las identidades de
suma de \'angulos,
\begin{equation}
  \label{eq:suma_angulos}
  \begin{aligned}
    \cos(\vph+\ps)&=\cos\vph\cos\ps-\sin\vph\sin\ps,\\
    \sin(\vph+\ps)&=\sin\vph\cos\ps+\sin\ps\cos\vph,
  \end{aligned}
\end{equation}
se obtiene
\begin{equation}
  \label{eq:z_cuadrado}
  \begin{aligned}
    z^2&=\abs{z}^2\qty{\cos\vph+i\sin\vph}^2\\
       &=\abs{z}^2\qty{\cos^2\vph-\sin^2\vph
          +2i\cos\vph\sin\vph}\\
       &=\abs{z}^2\qty{\cos(2\vph)+i\sin(2\vph)},
  \end{aligned}
\end{equation}
y multiplicando de nuevo por $z$,
\begin{equation}
  \label{eq:z_cubo}
  z^3=z^2z=\abs{z}^3\qty{\cos(3\vph)+i\sin(3\vph)} .
\end{equation}

\begin{teorema}[De Moivre]
  \label{teo:moivre}
  \begin{equation}
    \label{eq:moivre}
    z^n=\abs{z}^n\qty{\cos(n\vph)+i\sin(n\vph)}
    \qquad\forall n\in\N .
  \end{equation}
\end{teorema}

\begin{demostracion}[Demostraci\'on por inducci\'on]
  Se verifica para $n=1,2,3$ por \eqref{eq:polar_compleja},
  \eqref{eq:z_cuadrado} y \eqref{eq:z_cubo}. Suponi\'endolo v\'alido para
  $n-1$,
  \begin{equation}
    \label{eq:induccion_moivre}
    \small
    \begin{aligned}
      z^n=zz^{n-1}
        &=\abs{z}\qty{\cos\vph+i\sin\vph}\\
        &\quad\cdot\abs{z}^{n-1}
          \qty{\cos((n-1)\vph)+i\sin((n-1)\vph)}\\
        &=\abs{z}^n\qty{\cos(n\vph)+i\sin(n\vph)},
    \end{aligned}
  \end{equation}
  usando otra vez \eqref{eq:suma_angulos}.
\end{demostracion}

\begin{ejemplo}
  \label{ej:potencia_octava}
  Sea $z=1-i$, de modo que $\abs{z}=\sqrt{2}$ y
  $\Argu(z)=-\tfrac{\pi}{4}$. Entonces
  \begin{equation}
    \label{eq:potencia_octava}
    \begin{aligned}
      z^8&=\abs{z}^8\qty{\cos(8\vph)+i\sin(8\vph)}\\
         &=16\qty{\cos(-2\pi)+i\sin(-2\pi)}=16 .
    \end{aligned}
  \end{equation}
\end{ejemplo}

\subsection{Ra\'ices}
\label{subsec:raices}

Supongamos $z=\abs{z}\qty{\cos\vph+i\sin\vph}$ y busquemos $z^{1/m}$ con
$m\in\N$. Es equivalente a encontrar $w=\abs{w}\qty{\cos\ps+i\sin\ps}$ tal que
$z=w^m$. Usando \eqref{eq:moivre} en el miembro derecho y
\eqref{eq:argumento_multivaluado} en el izquierdo,
\begin{equation}
  \label{eq:raiz_igualdad}
  \small
  \abs{z}\qty{\cos(\vph+2n\pi)+i\sin(\vph+2n\pi)}
    =\abs{w}^m\qty{\cos(m\ps)+i\sin(m\ps)} .
\end{equation}
Identificando t\'erminos,
\begin{equation}
  \label{eq:raiz_identificacion}
  \abs{w}=\abs{z}^{1/m},
  \qquad
  m\ps=\vph+2n\pi
  \imp
  \ps=\frac{\vph}{m}+\frac{2n\pi}{m},
\end{equation}
y pedir que $\Argu(w)$ sea principal restringe a
$n=0,1,2,\dots,m-1$. Es decir, hay exactamente $m$ ra\'ices,
\begin{equation}
  \label{eq:raices_m}
  \small
  w_{\ell}=\abs{z}^{1/m}
    \qty{\cos\qty{\frac{\vph}{m}+\frac{2\ell\pi}{m}}
        +i\sin\qty{\frac{\vph}{m}+\frac{2\ell\pi}{m}}},
\end{equation}
con $\ell=0,1,\dots,m-1$, equiespaciadas sobre una circunferencia de radio
$\abs{z}^{1/m}$.

\begin{ejemplo}
  \label{ej:raiz_octava}
  Sea $z=16$, con $\vph=0$. Entonces $\abs{w}=16^{1/8}=\sqrt{2}$ y
  \begin{equation}
    \label{eq:raiz_octava}
    z^{1/8}=\sqrt{2}\qty{\cos\qty{\frac{n\pi}{4}}
      +i\sin\qty{\frac{n\pi}{4}}},
    \quad n=0,\dots,7,
  \end{equation}
  es decir
  \begin{equation}
    \label{eq:raices_16}
    \small
    16^{1/8}=\set{\sqrt{2},\,1+i,\,\sqrt{2}i,\,-1+i,\,
      -\sqrt{2},\,-1-i,\,-\sqrt{2}i,\,1-i} .
  \end{equation}
\end{ejemplo}

\begin{observacion}[Punto rama]
  Escribiendo $z=\abs{z}w$ con $w=\cos\vph+i\sin\vph$, al variar $\vph$ se
  recorre una circunferencia. La consecuencia fuerte de
  \eqref{eq:raices_m} es que en $z=0$ la ra\'iz no est\'a bien definida: todas
  las determinaciones colapsan. Decimos que $z=0$ es un \textbf{punto rama}.
\end{observacion}

\subsection{El plano complejo extendido}
\label{subsec:plano_extendido}

En el plano complejo no est\'a definido $\infty$, de modo que hay que
completarlo:
\begin{equation}
  \label{eq:plano_extendido}
  \bar{\C}=\C\cup\set{\infty},
\end{equation}
con las convenciones $a+\infty=\infty$ para todo $a\in\C$ y
\begin{equation}
  \label{eq:infinito}
  \lim_{\ep\rightarrow 0}\frac{a}{\ep}=\infty,
  \qquad a\neq 0 .
\end{equation}

\subsubsection{Proyecci\'on estereogr\'afica}
\label{subsubsec:proyeccion_estereografica}

El plano complejo es isomorfo a $\R^2$, lo que permite representarlo sobre una
esfera. Con $u,v,w\in\R$ sobre la esfera unitaria $u^2+v^2+w^2=1$ definimos
\begin{equation}
  \label{eq:proyeccion_estereografica}
  z=\frac{u+iv}{1-w} .
\end{equation}
El polo $w=1$ corresponde entonces al punto del infinito,
\begin{equation}
  \label{eq:polo_infinito}
  \lim_{w\rightarrow 1}z=\lim_{w\rightarrow 1}\frac{u+iv}{1-w}=\infty .
\end{equation}

\subsection{L\'imites y continuidad}
\label{subsec:limites_complejos}

\begin{definicion}[L\'imite]
  \label{def:limite_complejo}
  Sea $f(z)$ una funci\'on en el plano complejo, an\'aloga a una
  $\vb{f}:\R^2\rightarrow\R^2$. Decimos que
  \begin{equation}
    \label{eq:limite_complejo}
    \lim_{z\rightarrow z_0}f(z)=L
  \end{equation}
  si $\forall\ep>0$ existe $\de_\ep>0$ tal que
  $0<\abs{z-z_0}<\de_\ep$ implica $\abs{f(z)-L}<\ep$.
\end{definicion}

De $\abs{z}=\qty{z\bar{z}}^{1/2}=\abs{\bar{z}}$ se sigue de inmediato
\begin{equation}
  \label{eq:limite_conjugado}
  \lim_{z\rightarrow z_0}\bar{f}(z)=\bar{L},
\end{equation}
y como consecuencia
\begin{equation}
  \label{eq:limite_partes}
  \lim_{z\rightarrow z_0}\Real\set{f(z)}=\Real\set{L},
  \qquad
  \lim_{z\rightarrow z_0}\Imag\set{f(z)}=\Imag\set{L},
\end{equation}
donde, si $z=x+iy$,
\begin{equation}
  \label{eq:partes_real_imaginaria}
  \Real\set{z}=x=\frac{z+\bar{z}}{2},
  \qquad
  \Imag\set{z}=y=\frac{z-\bar{z}}{2i} .
\end{equation}

\begin{definicion}[Continuidad]
  \label{def:continuidad_compleja}
  $f(z)$ es continua en $z_0\in\Om$ si
  $\lim_{z\rightarrow z_0}f(z)=f(z_0)$.
\end{definicion}

\subsection{Derivada y ecuaciones de Cauchy--Riemann}
\label{subsec:cauchy_riemann}

\begin{definicion}[Derivada compleja]
  \label{def:derivada_compleja}
  $f(z)$ es derivable en $z_0\in\Om$ si existe
  \begin{equation}
    \label{eq:derivada_compleja}
    \dv{}{z}f(z)=\lim_{\al\rightarrow 0}
      \frac{f(z+\al)-f(z)}{\al},
    \qquad \al\in\C .
  \end{equation}
\end{definicion}

\begin{observacion}
  La condici\'on no s\'olo se pide en $z_0$ sino en toda una vecindad de
  $z_0$; \'esa es la diferencia esencial con el caso real, y de ella se sigue
  todo lo dem\'as.
\end{observacion}

Lo decisivo es que $\al\in\C$ puede tender a cero por cualquier direcci\'on.
Consideremos los dos acercamientos particulares
$\al=h$ (paralelo al eje real) y $\al=ih$ (paralelo al eje imaginario), con
$h\in\R$. Escribiendo
\begin{equation}
  \label{eq:f_u_v}
  f(z)=f(x,y)=u(x,y)+iv(x,y),
  \qquad u,v:\R^2\rightarrow\R,
\end{equation}
el primero da
\begin{equation}
  \label{eq:derivada_x}
  \dv{}{z}f(z)=\lim_{h\rightarrow 0}\frac{f(x+h,y)-f(x,y)}{h}
    =\pdv{}{x}f(z),
\end{equation}
y el segundo
\begin{equation}
  \label{eq:derivada_y}
  \dv{}{z}f(z)=\lim_{h\rightarrow 0}\frac{f(x,y+h)-f(x,y)}{ih}
    =-i\pdv{}{y}f(z).
\end{equation}
Como $f$ es derivable, ambos l\'imites deben coincidir:
\begin{equation}
  \label{eq:condicion_derivada}
  \dv{}{z}f(z)=\pdv{}{x}f(z)=-i\pdv{}{y}f(z).
\end{equation}

\begin{ejemplo}
  \label{ej:z_cuadrado_derivada}
  Sea $f(z)=z^2=(x+iy)^2=x^2-y^2+2ixy$. Entonces
  \begin{equation}
    \label{eq:z2_derivadas}
    \begin{aligned}
      \pdv{}{x}f(z)&=2x+2iy=2z,\\
      -i\pdv{}{y}f(z)&=-i\qty{-2y+2ix}=2x+2iy=2z,
    \end{aligned}
  \end{equation}
  y en efecto $\dv{}{z}z^2=2z$.
\end{ejemplo}

Separando \eqref{eq:condicion_derivada} en partes real e imaginaria,
\begin{equation}
  \label{eq:cr_desarrollo}
  \begin{aligned}
    \pdv{}{x}f&=\pdv{u}{x}+i\pdv{v}{x},\\
    -i\pdv{}{y}f&=\pdv{v}{y}-i\pdv{u}{y},
  \end{aligned}
\end{equation}
e igualando componente a componente obtenemos el resultado central.

\begin{teorema}[Ecuaciones de Cauchy--Riemann]
  \label{teo:cauchy_riemann}
  Si $f=u+iv$ es derivable, entonces
  \begin{equation}
    \label{eq:cauchy_riemann}
    \pdv{u}{x}=\pdv{v}{y},
    \qquad
    \pdv{u}{y}=-\pdv{v}{x} .
  \end{equation}
\end{teorema}

\begin{observacion}[Relaci\'on con el jacobiano]
  \label{obs:jacobiano_cr}
  Usando \eqref{eq:cauchy_riemann},
  \begin{equation}
    \label{eq:modulo_derivada}
    \begin{aligned}
      \abs{\dv{}{z}f(z)}^2
        &=\qty{\pdv{u}{x}}^2+\qty{\pdv{v}{x}}^2\\
        &=\pdv{u}{x}\pdv{v}{y}-\pdv{v}{x}\pdv{u}{y}\\
        &=
        \begin{vmatrix}
          \pdv{u}{x} & \pdv{v}{x}\\[3pt]
          \pdv{u}{y} & \pdv{v}{y}
        \end{vmatrix}
        =\abs{\frac{\partial(u,v)}{\partial(x,y)}} ,
    \end{aligned}
  \end{equation}
  el jacobiano de la transformaci\'on. Por el teorema de la funci\'on inversa
  (Teorema~\ref{teo:funcion_inversa}), si $\dv{}{z}f(z)\neq 0$ existe
  localmente la transformaci\'on inversa.
\end{observacion}

\subsection{Funciones arm\'onicas}
\label{subsec:armonicas}

Si $f(z)$ es diferenciable, sus derivadas sucesivas son continuas. Entonces,
aplicando dos veces \eqref{eq:cauchy_riemann},
\begin{equation}
  \label{eq:laplace_u}
  \begin{aligned}
    \pddv{u}{x}+\pddv{u}{y}
      &=\pdv{}{x}\qty{\pdv{u}{x}}+\pdv{}{y}\qty{\pdv{u}{y}}\\
      &=\pdv{}{x}\qty{\pdv{v}{y}}+\pdv{}{y}\qty{-\pdv{v}{x}}=0 ,
  \end{aligned}
\end{equation}
y an\'alogamente para $v$. Es decir,
\begin{equation}
  \label{eq:laplace}
  \Lap u=0,
  \qquad
  \Lap v=0,
  \qquad
  \Lap=\pddv{}{x}+\pddv{}{y},
\end{equation}
la \textbf{ecuaci\'on de Laplace}: las partes real e imaginaria de una
funci\'on anal\'itica son \textbf{arm\'onicas}.

\begin{observacion}
  Las funciones arm\'onicas alcanzan su m\'aximo y su m\'inimo en la frontera
  del dominio.
\end{observacion}

\subsection{Polinomios}
\label{subsec:polinomios}

Sea $f(z)=P_m(z)=a_0+a_1z+\dots+a_mz^m$ con $a_m\neq 0$. De
\eqref{eq:condicion_derivada} se sigue
\begin{equation}
  \label{eq:derivada_potencia_compleja}
  \dv{}{z}z^k=\pdv{}{x}z^k=-i\pdv{}{y}z^k=kz^{k-1} .
\end{equation}
Derivando repetidamente y evaluando en $z=0$,
\begin{equation}
  \label{eq:coeficientes_polinomio}
  \eval{\dnv{}{z}{k}P_m(z)}_{z=0}=k!\,a_k
  \imp
  a_k=\frac{1}{k!}\eval{\dnv{}{z}{k}P_m(z)}_{z=0},
\end{equation}
para $k\le m$.

Si $\al_1$ es una ra\'iz, $P_m(\al_1)=0$, entonces
$P_m(z)=(z-\al_1)P_{m-1}(z)$ con $P_{m-1}(\al_1)\neq 0$. Iterando,
\begin{equation}
  \label{eq:factorizacion}
  P_m(z)=a_m(z-\al_1)(z-\al_2)\dots(z-\al_m),
\end{equation}
y desarrollando alrededor de un punto $z_0$ cualquiera,
\begin{equation}
  \label{eq:polinomio_centrado}
  P_m(z)=b_0+b_1(z-z_0)+b_2(z-z_0)^2+\dots+b_m(z-z_0)^m .
\end{equation}

\section{Series de potencias}
\label{sec:series_potencia}

\subsection{Sucesiones}
\label{subsec:sucesiones}

Consideremos una sucesi\'on de n\'umeros $\set{a_n}_{n=1}^{\infty}$. Decimos
que $A$ es su l\'imite si $\forall\ep>0$ existe $n_0$ tal que
$\abs{a_n-A}<\ep$ para todo $n\ge n_0$.

\begin{definicion}[Convergencia]
  \label{def:convergencia}
  Una sucesi\'on con l\'imite finito se dice \textbf{convergente}; si no
  converge es \textbf{divergente}.
\end{definicion}

\begin{definicion}[Sucesi\'on de Cauchy]
  \label{def:cauchy}
  Una sucesi\'on se llama \textbf{fundamental} o \textbf{de Cauchy} si
  $\forall\ep>0$ existe $n_0$ tal que
  \begin{equation}
    \label{eq:cauchy}
    \abs{a_n-a_m}<\ep
    \qquad\forall n,m\ge n_0 .
  \end{equation}
\end{definicion}

\begin{ejemplo}
  \label{ej:sucesion_cauchy}
  Sea $a_n=1-\dfrac{3}{n+2}$. Entonces
  \begin{equation}
    \label{eq:ejemplo_cauchy}
    \begin{aligned}
      \abs{a_n-a_m}
        &=\abs{1-\frac{3}{n+2}-\qty{1-\frac{3}{m+2}}}\\
        &=3\abs{\frac{n-m}{(n+2)(m+2)}}<\ep
    \end{aligned}
  \end{equation}
  para $n,m$ suficientemente grandes, de modo que la sucesi\'on es de Cauchy.
\end{ejemplo}

\begin{observacion}
  Las partes real e imaginaria de una sucesi\'on de Cauchy son sucesiones de
  Cauchy.
\end{observacion}

\subsection{Series}
\label{subsec:series}

Sea la serie infinita
\begin{equation}
  \label{eq:serie}
  S=a_0+a_1+a_2+\dots+a_n+\dots
\end{equation}
con $\set{a_n}$ los elementos de una sucesi\'on. Definimos las
\textbf{sumas parciales} $S_n=a_0+a_1+\dots+a_n$ y con ellas una nueva
sucesi\'on $\set{S_n}_{n=1}^{\infty}$.

Una serie infinita \textbf{converge} si $\forall\ep>0$ existe $n_0$ tal que
$\abs{S_p-S_n}<\ep$ para todos $n,p\ge n_0$; es decir, si la sucesi\'on de
sumas parciales es de Cauchy. En ese caso podemos hablar del
\textbf{residuo}:
\begin{equation}
  \label{eq:residuo}
  S=S_n+R_n,
  \qquad
  R_n=\sum_{j=n+1}^{\infty}a_j .
\end{equation}

Comparando $S$ con la serie de valores absolutos
\begin{equation}
  \label{eq:serie_absoluta}
  S_a=\abs{a_1}+\abs{a_2}+\dots+\abs{a_n}+\dots
\end{equation}
y usando la desigualdad triangular,
\begin{equation}
  \label{eq:desigualdad_triangular_serie}
  \abs{a_n+a_{n+1}+\dots+a_p}
    \le\abs{a_n}+\abs{a_{n+1}}+\dots+\abs{a_p},
\end{equation}
se ve que si $S_a$ converge entonces $S$ converge. En tal caso se dice que la
serie es \textbf{absolutamente convergente}.

\subsection{Convergencia uniforme}
\label{subsec:convergencia_uniforme}

Consideremos ahora una sucesi\'on de funciones definidas en un dominio $\Om$.
Si $\set{f_n(z)}_{n=1}^{\infty}$ converge para todo $z\in\Om$, el l\'imite
$f(z)$ es una funci\'on definida en $\Om$: dados $\ep>0$ y $z\in\Om$, existe
$n_0$ tal que $\abs{f_n(z)-f(z)}<\ep$ para $n\ge n_0$. Pero ese $n_0$ depende
en general de $z$.

\begin{definicion}[Convergencia uniforme]
  \label{def:convergencia_uniforme}
  Una sucesi\'on $\set{f_n(z)}$ converge \textbf{uniformemente} a $f(z)$ en
  $\Om$ si $\forall\ep>0$ existe un $n_0$ \emph{independiente de $z$} tal que
  \begin{equation}
    \label{eq:convergencia_uniforme}
    \abs{f_n(z)-f_m(z)}<\ep
    \qquad\forall n,m\ge n_0 .
  \end{equation}
\end{definicion}

\begin{observacion}
  La funci\'on l\'imite de una sucesi\'on uniformemente convergente de
  funciones continuas es continua.
\end{observacion}

\subsection{Caso de estudio: la serie geom\'etrica}
\label{subsec:serie_geometrica}

Sea
\begin{equation}
  \label{eq:geometrica}
  S=1+z+z^2+\dots+z^n+\dots,
  \qquad
  S_n=1+z+z^2+\dots+z^n .
\end{equation}

\begin{observacion}[El truco]
  Multiplicando por $z$ y restando,
  \begin{equation}
    \label{eq:truco_geometrica}
    \begin{aligned}
      zS_n&=z+z^2+\dots+z^{n+1}\\
      \imp S_n-zS_n&=1-z^{n+1}\\
      \imp S_n&=\frac{1-z^{n+1}}{1-z},\qquad z\neq 1 .
    \end{aligned}
  \end{equation}
\end{observacion}

Entonces
\begin{equation}
  \label{eq:limite_geometrica}
  S=\lim_{n\rightarrow\infty}S_n
   =\lim_{n\rightarrow\infty}\frac{1-z^{n+1}}{1-z},
\end{equation}
y se distinguen dos casos:
\begin{enumerate}
  \item si $\abs{z}<1$, entonces $z^{n+1}\rightarrow 0$ y
        $S=\dfrac{1}{1-z}$;
  \item si $\abs{z}>1$, entonces $\abs{z^{n+1}}\rightarrow\infty$ y la serie
        diverge.
\end{enumerate}

\begin{observacion}
  Existe pues un valor $R$ tal que hay convergencia si $\abs{z}<R$. Esto se
  generaliza a cualquier serie de potencias.
\end{observacion}

\subsection{Radio de convergencia}
\label{subsec:radio_convergencia}

\begin{teorema}[Radio de convergencia]
  \label{teo:radio_convergencia}
  Para cada serie de potencias existe un valor $R$, llamado \textbf{radio de
  convergencia}, con las siguientes propiedades:
  \begin{itemize}
    \item la serie converge absolutamente para $\abs{z}<R$, y si
          $0\le\rho<R$ la convergencia es uniforme para $\abs{z}\le\rho$;
    \item si $\abs{z}>R$ los t\'erminos de la serie no tienen cota y la serie
          diverge;
    \item en $\abs{z}<R$ la suma de la serie es una funci\'on anal\'itica; su
          derivada se obtiene diferenciando t\'ermino a t\'ermino y la serie
          de las derivadas tiene el mismo radio de convergencia.
  \end{itemize}
  Adem\'as vale la \textbf{f\'ormula de Hadamard},
  \begin{equation}
    \label{eq:hadamard}
    \frac{1}{R}=\limsup_{n\rightarrow\infty}\;\abs{a_n}^{1/n} .
  \end{equation}
\end{teorema}

\subsection{Teorema de Taylor}
\label{subsec:taylor_complejo}

\begin{teorema}[Taylor]
  \label{teo:taylor}
  Si $f(z)$ es anal\'itica en $\Om$, entonces
  \begin{equation}
    \label{eq:taylor_complejo}
    f(z)=\sum_{k=0}^{\infty}\frac{1}{k!}
      \eval{\dnv{}{z}{k}f(z)}_{z=z_0}(z-z_0)^k .
  \end{equation}
\end{teorema}

\subsection{Funciones anal\'iticas elementales}
\label{subsec:funciones_analiticas}

\subsubsection{Seno y coseno}
\label{subsubsec:seno_coseno}

Desarrollando $\sin z$ alrededor de $z=0$ con \eqref{eq:taylor_complejo},
\begin{equation}
  \label{eq:serie_seno}
  \sin z=z-\frac{1}{3!}z^3+\frac{1}{5!}z^5-\dots
    =\sum_{k=0}^{\infty}\frac{(-1)^k}{(2k+1)!}z^{2k+1},
\end{equation}
y derivando t\'ermino a t\'ermino, lo cual es l\'icito por el
Teorema~\ref{teo:radio_convergencia},
\begin{equation}
  \label{eq:serie_coseno}
  \cos z=\dv{}{z}\sin z
    =\sum_{k=0}^{\infty}\frac{(-1)^k}{(2k)!}z^{2k} .
\end{equation}

\begin{observacion}[Radio de convergencia]
  Para \eqref{eq:serie_seno} los coeficientes son
  $a_k=(-1)^k/(2k+1)!$. Como $k^k<(2k+1)!$ para $k$ grande,
  \begin{equation}
    \label{eq:cota_factorial}
    \frac{1}{(2k+1)!}<\frac{1}{k^k}
    \imp
    \qty{\frac{1}{(2k+1)!}}^{1/k}<\frac{1}{k},
  \end{equation}
  y como $\lim_{k\rightarrow\infty}1/k=0$, la f\'ormula de Hadamard
  \eqref{eq:hadamard} da $1/R=0$: la serie de $\sin z$ tiene radio de
  convergencia infinito.
\end{observacion}

\subsubsection{Relaci\'on con las funciones hiperb\'olicas}
\label{subsubsec:hiperbolicas}

Recordemos
\begin{equation}
  \label{eq:hiperbolicas}
  \sinh\al=\frac{e^{\al}-e^{-\al}}{2},
  \qquad
  \cosh\al=\frac{e^{\al}+e^{-\al}}{2},
\end{equation}
que satisfacen $\dv{}{\al}\sinh\al=\cosh\al$ y
$\dv{}{\al}\cosh\al=\sinh\al$. Desarrollando $\sinh$ alrededor de cero,
\begin{equation}
  \label{eq:serie_sinh}
  \sinh\al=\sum_{k=0}^{\infty}\frac{\al^{2k+1}}{(2k+1)!} .
\end{equation}
Por otro lado, evaluando \eqref{eq:serie_seno} en $z=i\al$ y usando
$i^{2k+1}=i(-1)^k$,
\begin{equation}
  \label{eq:seno_imaginario}
  \sin(i\al)=\sum_{k=0}^{\infty}\frac{(-1)^k}{(2k+1)!}(i\al)^{2k+1}
    =i\sum_{k=0}^{\infty}\frac{\al^{2k+1}}{(2k+1)!} ,
\end{equation}
de modo que, comparando con \eqref{eq:serie_sinh},
\begin{equation}
  \label{eq:seno_sinh}
  \sin(i\al)=i\sinh\al .
\end{equation}
An\'alogamente,
\begin{equation}
  \label{eq:coseno_cosh}
  \cos(i\al)=\sum_{k=0}^{\infty}\frac{\al^{2k}}{(2k)!}=\cosh\al .
\end{equation}

\subsubsection{La exponencial}
\label{subsubsec:exponencial}

Invirtiendo \eqref{eq:hiperbolicas} con $\al=iy$ se obtienen
\begin{equation}
  \label{eq:seno_coseno_exponencial}
  \sin y=\frac{e^{iy}-e^{-iy}}{2i},
  \qquad
  \cos y=\frac{e^{iy}+e^{-iy}}{2},
\end{equation}
y sum\'andolas, la \textbf{f\'ormula de Euler}
\begin{equation}
  \label{eq:euler}
  e^{iy}=\cos y+i\sin y .
\end{equation}

\begin{definicion}[Exponencial compleja]
  \label{def:exponencial}
  La funci\'on $e^z$ con $z\in\C$ es anal\'itica y vale su serie de Taylor,
  \begin{equation}
    \label{eq:serie_exponencial}
    e^z=\sum_{k=0}^{\infty}\frac{1}{k!}z^k .
  \end{equation}
\end{definicion}

Su radio de convergencia tambi\'en es infinito: como $k^{k/10}<k!$ para $k$
suficientemente grande,
\begin{equation}
  \label{eq:radio_exponencial}
  \qty{\frac{1}{k!}}^{1/k}<\qty{\frac{1}{k^{k/10}}}^{1/k}
    =\frac{1}{k^{1/10}}\xrightarrow[k\rightarrow\infty]{}0 .
\end{equation}

Separando \eqref{eq:serie_exponencial} para $z=iy$ en las potencias pares e
impares se recupera \eqref{eq:euler}:
\begin{equation}
  \label{eq:exponencial_euler}
  \small
  e^{iy}=\sum_{\ell=0}^{\infty}\frac{(-1)^\ell}{(2\ell)!}y^{2\ell}
    +i\sum_{\ell=0}^{\infty}\frac{(-1)^\ell}{(2\ell+1)!}y^{2\ell+1}
    =\cos y+i\sin y .
\end{equation}

En consecuencia, la representaci\'on polar \eqref{eq:polar_compleja} se
escribe de forma compacta
\begin{equation}
  \label{eq:polar_exponencial}
  z=\abs{z}\qty{\cos\vph+i\sin\vph}=\abs{z}e^{i\vph},
  \qquad \vph=\Argu(z),
\end{equation}
y para $z=x+iy$,
\begin{equation}
  \label{eq:exponencial_cartesiana}
  e^{z}=e^{x+iy}=e^{x}\qty{\cos y+i\sin y} .
\end{equation}

\begin{ejemplo}
  \label{ej:polar_exponencial}
  \begin{equation}
    \label{eq:ejemplos_polar}
    z_1=1+i=\sqrt{2}e^{i\pi/4},
    \qquad
    z_2=i=e^{i\pi/2} .
  \end{equation}
\end{ejemplo}

\subsubsection{El logaritmo}
\label{subsubsec:logaritmo}

\begin{definicion}[Logaritmo complejo]
  \label{def:logaritmo}
  \begin{equation}
    \label{eq:logaritmo}
    \ln(z)=\ln\qty{\abs{z}e^{i\vph}}=\ln\abs{z}+i\vph,
    \qquad \vph=\Argu(z).
  \end{equation}
\end{definicion}

De \eqref{eq:logaritmo} se siguen las propiedades usuales. Para el producto,
\begin{equation}
  \label{eq:log_producto}
  \begin{aligned}
    \ln(z_1z_2)
      &=\ln\qty{\abs{z_1z_2}e^{i(\vph_1+\vph_2)}}\\
      &=\ln\abs{z_1}+i\vph_1+\ln\abs{z_2}+i\vph_2\\
      &=\ln(z_1)+\ln(z_2),
  \end{aligned}
\end{equation}
y para el cociente, con $z_2\neq 0$,
\begin{equation}
  \label{eq:log_cociente}
  \begin{aligned}
    \ln\qty{\frac{z_1}{z_2}}
      &=\ln\qty{\frac{z_1\bar{z}_2}{\abs{z_2}^2}}
       =\ln(z_1)+\ln\qty{\frac{\bar{z}_2}{\abs{z_2}^2}}\\
      &=\ln(z_1)+\ln\qty{\frac{1}{\abs{z_2}}}-i\vph_2\\
      &=\ln(z_1)-\ln(z_2).
  \end{aligned}
\end{equation}

\begin{observacion}
  Como el argumento est\'a definido salvo m\'ultiplos de $2\pi$, el logaritmo
  complejo es multivaluado y $z=0$ vuelve a ser un punto rama.
\end{observacion}

\section{Series de Fourier}
\label{sec:fourier}

Supongamos que $f$ es continua por pedazos en $x\in[a,b]$ y suave, es decir,
que $\dv{}{x}f(x)$ es continua por pedazos en $(a,b)$, y pidamos adem\'as que
existan $\eval{\dv{}{x}f(x)}_{a^{+}}$ y $\eval{\dv{}{x}f(x)}_{b^{-}}$. De ese
universo de funciones nos interesan las peri\'odicas de periodo $2\pi$.

Si $f$ admite una expansi\'on en alguna base $\set{g_k}$ como en
\eqref{eq:expansion_general}, $f(x)=\sum_k C_kg_k(x)$, la elecci\'on
\begin{equation}
  \label{eq:base_fourier}
  g_k(x)=\qty{e^{ix}}^{k}=e^{ikx}
\end{equation}
produce la \textbf{serie de Fourier}
\begin{equation}
  \label{eq:serie_fourier}
  f(x)=\sum_{k=-\infty}^{\infty}C_ke^{ikx} .
\end{equation}

\begin{teorema}[Convergencia]
  \label{teo:convergencia_fourier}
  Si $f$ es peri\'odica de periodo $2\pi$, continua y suave por pedazos,
  entonces su serie de Fourier converge absoluta y uniformemente a $f$.
\end{teorema}

\subsection{Par\'entesis: producto interno de funciones}
\label{subsec:producto_interno_funciones}

Para vectores, el producto interno es
$\braket{\vb{u}}{\vb{v}}=\vb{u}^\dg\vb{v}$. Para funciones se define de
manera an\'aloga sustituyendo la suma sobre componentes por una integral,
\begin{equation}
  \label{eq:pi_funciones}
  \braket{g}{f}=\frac{1}{N}\int_{a}^{b}g^{*}(x)f(x)\dd x,
\end{equation}
donde $N$ es una constante de normalizaci\'on. Apliqu\'emoslo a la base
\eqref{eq:base_fourier}:
\begin{equation}
  \label{eq:pi_exponenciales}
  \braket{e^{ik'x}}{e^{ikx}}
    =\frac{1}{N}\int_{-\pi}^{\pi}e^{-ik'x}e^{ikx}\dd x .
\end{equation}
La integral se calcula directamente,
\begin{equation}
  \label{eq:integral_ortogonalidad}
  \int_{-\pi}^{\pi}e^{i(k-k')x}\dd x=
  \begin{cases}
    2\pi, & k=k',\\
    0, & k\neq k',
  \end{cases}
  =2\pi\de_{kk'},
\end{equation}
de modo que la normalizaci\'on adecuada es $N=2\pi$ y
\begin{equation}
  \label{eq:ortonormalidad_fourier}
  \braket{e^{ik'x}}{e^{ikx}}
    =\frac{1}{2\pi}\int_{-\pi}^{\pi}e^{-ik'x}e^{ikx}\dd x=\de_{kk'} .
\end{equation}
Es decir, respecto del producto interno
\begin{equation}
  \label{eq:pi_fourier}
  \braket{f}{g}=\frac{1}{2\pi}\int_{-\pi}^{\pi}f^{*}(x)g(x)\dd x,
\end{equation}
la familia $\set{e^{ikx}}_{k=-\infty}^{\infty}$ es una \textbf{base
ortonormal}.

\subsection{Los coeficientes}
\label{subsec:coeficientes_fourier}

La ortonormalidad permite extraer los coeficientes proyectando, exactamente
como en \eqref{eq:contraccion_kronecker}:
\begin{equation}
  \label{eq:extraccion_coeficientes}
  \begin{aligned}
    \braket{e^{ik'x}}{f(x)}
      &=\frac{1}{2\pi}\int_{-\pi}^{\pi}e^{-ik'x}
        \sum_{k=-\infty}^{\infty}C_ke^{ikx}\dd x\\
      &=\sum_{k=-\infty}^{\infty}C_k\,
        \frac{1}{2\pi}\int_{-\pi}^{\pi}e^{-ik'x}e^{ikx}\dd x\\
      &=\sum_{k=-\infty}^{\infty}C_k\de_{kk'}=C_{k'} ,
  \end{aligned}
\end{equation}
es decir
\begin{equation}
  \label{eq:coeficientes_fourier}
  C_k=\braket{e^{ikx}}{f(x)}
    =\frac{1}{2\pi}\int_{-\pi}^{\pi}e^{-ikx}f(x)\dd x .
\end{equation}

\subsection{Forma real}
\label{subsec:fourier_real}

Agrupando en \eqref{eq:serie_fourier} los t\'erminos $\pm k$ y usando la
f\'ormula de Euler \eqref{eq:euler},
\begin{equation}
  \label{eq:fourier_real_desarrollo}
  \small
  \begin{aligned}
    f(x)&=\dots+C_{-1}e^{-ix}+C_0+C_1e^{ix}+\dots\\
        &=C_0+\dots+C_{-1}\qty{\cos x-i\sin x}\\
        &\qquad\quad+C_{1}\qty{\cos x+i\sin x}+\dots\\
        &=C_0+\dots+(C_{1}+C_{-1})\cos x\\
        &\qquad\quad+i(C_{1}-C_{-1})\sin x+\dots
  \end{aligned}
\end{equation}
Definiendo $a_k=C_k+C_{-k}$ y $b_k=i(C_k-C_{-k})$ se obtiene la
\textbf{forma real} de la serie,
\begin{equation}
  \label{eq:fourier_real}
  f(x)=\frac{a_0}{2}+\sum_{k=1}^{\infty}a_k\cos(kx)
    +\sum_{k=1}^{\infty}b_k\sin(kx).
\end{equation}

Las familias $\set{\cos(kx)}_{k=0}^{\infty}$ y
$\set{\sin(kx)}_{k=1}^{\infty}$ son tambi\'en ortonormales, ahora respecto de
la normalizaci\'on $N=\pi$. En efecto, usando la identidad producto-suma,
\begin{equation}
  \label{eq:ortonormalidad_senos}
  \small
  \begin{aligned}
    \braket{\sin(k'x)}{\sin(kx)}
      &=\frac{1}{\pi}\int_{-\pi}^{\pi}\sin(k'x)\sin(kx)\dd x\\
      &=\frac{1}{2\pi}\int_{-\pi}^{\pi}
        \big(\cos((k-k')x)-\cos((k+k')x)\big)\dd x\\
      &=\de_{kk'},
  \end{aligned}
\end{equation}
y an\'alogamente
\begin{equation}
  \label{eq:ortonormalidad_cosenos}
  \braket{\cos(k'x)}{\cos(kx)}=\de_{kk'},
  \qquad
  \braket{\cos(k'x)}{\sin(kx)}=0
\end{equation}
para todos $k,k'\ge 1$. Por lo tanto, proyectando igual que en
\eqref{eq:extraccion_coeficientes},
\begin{equation}
  \label{eq:coeficientes_reales}
  \begin{aligned}
    a_k&=\braket{\cos(kx)}{f(x)}
        =\frac{1}{\pi}\int_{-\pi}^{\pi}f(x)\cos(kx)\dd x,\\
    b_k&=\braket{\sin(kx)}{f(x)}
        =\frac{1}{\pi}\int_{-\pi}^{\pi}f(x)\sin(kx)\dd x,\\
    \frac{a_0}{2}&=\frac{1}{2\pi}\int_{-\pi}^{\pi}f(x)\dd x ,
  \end{aligned}
\end{equation}
donde el \'ultimo t\'ermino es simplemente el valor medio de $f$ sobre el
periodo.

\section{Integraci\'on compleja}
\label{sec:integracion_compleja}

La integral \eqref{eq:integral_ortogonalidad} se separa en parte real e
imaginaria,
\begin{equation}
  \label{eq:separacion_integral}
  \small
  \int_{-\pi}^{\pi}e^{i(k-k')x}\dd x
    =\int_{-\pi}^{\pi}\cos((k-k')x)\dd x
    +i\int_{-\pi}^{\pi}\sin((k-k')x)\dd x,
\end{equation}
lo que da la pauta para definir la integral en el plano complejo.

\begin{definicion}[Integral compleja]
  \label{def:integral_compleja}
  Sea $f(z)$ bien definida en $D\subset\C$ y $\ga:\R\rightarrow\C$ una curva
  parametrizada contenida en $D$. Definimos
  \begin{equation}
    \label{eq:integral_compleja}
    \int_{\ga}f(z)\dd z=\int_{a}^{b}f(z(t))\dv{z}{t}\dd t .
  \end{equation}
\end{definicion}

La integral compleja es entonces una integral de l\'inea en el plano. Si
$z=x+iy$,
\begin{equation}
  \label{eq:dz}
  \dd z=\dv{z}{t}\dd t=\dv{}{t}(x+iy)\dd t
    =(\dot{x}+i\dot{y})\dd t=\dd x+i\dd y,
\end{equation}
de modo que la forma m\'as general posible es
\begin{equation}
  \label{eq:forma_general_integral}
  \int_{\ga}\qty{p\dd x+q\dd y},
  \qquad p,q\in\C .
\end{equation}

\subsection{Acotamiento}
\label{subsec:acotamiento_integral}

Se cumple
\begin{equation}
  \label{eq:cota_integral}
  \abs{\int_{\ga}f(z)\dd z}
    \le\int_{\ga}\abs{f(z)}\abs{\dv{z}{t}}\abs{\dd t} .
\end{equation}
Notando que si $z(t)=x(t)+iy(t)$,
\begin{equation}
  \label{eq:diferencial_arco_complejo}
  \abs{\dv{z}{t}}
    =\qty{\qty{\dv{x}{t}}^2+\qty{\dv{y}{t}}^2}^{1/2},
\end{equation}
que es el diferencial de longitud de arco, resulta en particular
\begin{equation}
  \label{eq:cota_longitud}
  \abs{\int_{\ga}\dd z}\le\int_{\ga}\abs{\dd z}
    =\text{longitud de arco de }\ga .
\end{equation}
Si $\ga$ es una circunferencia de radio $\rho$, entonces
\begin{equation}
  \label{eq:cota_circulo}
  \abs{\int_{\ga}\dd z}\le 2\pi\rho .
\end{equation}

\subsection{Independencia de la trayectoria}
\label{subsec:independencia_trayectoria}

Retomando \eqref{eq:forma_general_integral} con $p(z)$ y $q(z)$ complejas,
supongamos que existe $U(z)$ tal que
\begin{equation}
  \label{eq:potencial_U}
  \pdv{U}{x}=p,
  \qquad
  \pdv{U}{y}=q .
\end{equation}
Entonces, por la regla de la cadena,
\begin{equation}
  \label{eq:integral_exacta}
  \begin{aligned}
    \int_{\ga}\qty{p\dd x+q\dd y}
      &=\int_{a}^{b}\qty{\pdv{U}{x}\dv{x}{t}
        +\pdv{U}{y}\dv{y}{t}}\dd t\\
      &=\int_{a}^{b}\dv{}{t}U\dd t=U(b)-U(a).
  \end{aligned}
\end{equation}

\begin{teorema}
  \label{teo:independencia}
  La integral de l\'inea $\int_{\ga}\qty{p\dd x+q\dd y}$ definida sobre un
  dominio $D$ depende s\'olo de los puntos extremos de $\ga$ si y s\'olo si
  existe una funci\'on $U(z)=U(x,y)$ en $D$ con derivadas parciales
  \eqref{eq:potencial_U}.
\end{teorema}

\subsubsection{Consecuencias}
\label{subsubsec:consecuencias_independencia}

Tenemos una diferencial exacta si
\begin{equation}
  \label{eq:diferencial_exacta}
  f(z)\dd z=f(z)\dd x+if(z)\dd y
\end{equation}
proviene de una $F(z)$ en $D$ con
\begin{equation}
  \label{eq:F_parciales}
  \pdv{F}{x}=f(z),
  \qquad
  -i\pdv{F}{y}=f(z),
\end{equation}
es decir, con $\pdv{F}{x}=-i\pdv{F}{y}$, que por
\eqref{eq:condicion_derivada} es justamente la condici\'on de
Cauchy--Riemann. En otras palabras, $F(z)$ debe ser anal\'itica en $D$.

\begin{observacion}[Consecuencia fuerte]
  La integral de una funci\'on continua depende \'unicamente de los extremos
  de $\ga$ si y s\'olo si $f$ es la derivada de una funci\'on anal\'itica en
  $D$. En particular, si $\ga$ es cerrada,
  \begin{equation}
    \label{eq:integral_cerrada_nula}
    \oint_{\ga}f(z)\dd z=0
  \end{equation}
  siempre que $f(z)$ sea anal\'itica en el dominio encerrado por $\ga$. \'Este
  es el \textbf{teorema de Cauchy}.
\end{observacion}

\begin{ejemplo}
  \label{ej:potencias_cerrada}
  \begin{equation}
    \label{eq:potencias_cerrada}
    \oint_{\ga}(z-a)^n\dd z=0
    \qquad\forall n\ge 0,
  \end{equation}
  ya que la primitiva $F(z)=\dfrac{(z-a)^{n+1}}{n+1}$ es anal\'itica.
\end{ejemplo}

\begin{ejemplo}[El caso $n=-1$]
  \label{ej:integral_fundamental}
  Sea $\ga$ la circunferencia de radio $\rho$ centrada en $a$, parametrizada
  por
  \begin{equation}
    \label{eq:parametrizacion_circulo}
    z(\vph)=a+\rho e^{i\vph}
      =a+\rho\qty{\cos\vph+i\sin\vph},
    \quad\vph\in[0,2\pi],
  \end{equation}
  de modo que $\dd z=i\rho e^{i\vph}\dd\vph$. Entonces
  \begin{equation}
    \label{eq:integral_fundamental}
    \oint_{\abs{z-a}=\rho}\frac{\dd z}{z-a}
      =\int_{0}^{2\pi}\frac{i\rho e^{i\vph}}{\rho e^{i\vph}}\dd\vph
      =2\pi i ,
  \end{equation}
  resultado independiente de $\rho$ y central en todo lo que sigue.
\end{ejemplo}

\begin{ejemplo}
  \label{ej:integral_recta}
  Sea $\ga(t)=(1+i)t$ con $t\in[0,1]$, de modo que sobre $\ga$ se tiene
  $x=\Real\set{z}=t$ y $\dd z=(1+i)\dd t$. Entonces
  \begin{equation}
    \label{eq:integral_recta}
    \int_{\ga}x\dd z=\int_{0}^{1}t\,(1+i)\dd t
      =\eval{\frac{1}{2}(1+i)t^2}_{0}^{1}=\frac{1}{2}(1+i).
  \end{equation}
  N\'otese que el integrando $x$ no es anal\'itico, de modo que aqu\'i no
  aplica \eqref{eq:integral_cerrada_nula}.
\end{ejemplo}

\subsection{Deformaci\'on de contornos}
\label{subsec:deformacion_contornos}

Supongamos que $f(z)$ es anal\'itica en $D$ excepto en un punto $z=a$.
Consideremos dos curvas: $C_1$, deformable hasta $\ga$, y $C_2$, contenida en
el interior de $\ga$ y que encierra a $z=a$. Unimos ambas por un corte de
anchura $\ep$ para formar un contorno $C_3(\ep)$ que \emph{no} encierra a
$z=a$; entonces $f$ es anal\'itica sobre y dentro de $C_3$ y
\begin{equation}
  \label{eq:contorno_C3}
  \oint_{C_3}f(z)\dd z=0 .
\end{equation}
Al hacer $\ep\rightarrow 0$ las contribuciones del corte se cancelan y queda
\begin{equation}
  \label{eq:deformacion}
  \oint_{C_1}f(z)\dd z+\oint_{-C_2}f(z)\dd z=0
  \imp
  \oint_{\ga}f(z)\dd z=\oint_{C_2}f(z)\dd z .
\end{equation}
En particular podemos pedir que $C_2$ sea una circunferencia de radio $\rho$
centrada en la singularidad.

\begin{ejemplo}
  \label{ej:fracciones_parciales}
  \begin{equation}
    \label{eq:fracciones_parciales}
    \small
    \begin{aligned}
      \oint_{\abs{z}=2}\frac{\dd z}{z^2-1}
        &=\frac{1}{2}\oint_{\abs{z}=2}\frac{\dd z}{z-1}
         -\frac{1}{2}\oint_{\abs{z}=2}\frac{\dd z}{z+1}\\
        &=\frac{1}{2}\oint_{\abs{z-1}=\rho}\frac{\dd z}{z-1}
         -\frac{1}{2}\oint_{\abs{z+1}=\rho}\frac{\dd z}{z+1}\\
        &=\frac{1}{2}(2\pi i)-\frac{1}{2}(2\pi i)=0 ,
    \end{aligned}
  \end{equation}
  donde en el segundo paso se deform\'o cada contorno a un c\'irculo peque\~no
  alrededor de su polo y en el tercero se us\'o
  \eqref{eq:integral_fundamental}.
\end{ejemplo}

Generalizando: si $f(z)$ es anal\'itica en $D$ salvo en un conjunto finito de
puntos aislados $\set{z_j}_{j=1}^{n}$, entonces
\begin{equation}
  \label{eq:deformacion_general}
  \oint_{\ga}f(z)\dd z
    =\sum_{j=1}^{n}\oint_{B_\rho(z_j)}f(z)\dd z .
\end{equation}

\subsection{F\'ormula integral de Cauchy}
\label{subsec:formula_cauchy}

Sea $f(z)$ anal\'itica en $D$ y construyamos
\begin{equation}
  \label{eq:F_cociente}
  F(z)=\frac{f(z)-f(a)}{z-a} .
\end{equation}
Al ser $f$ anal\'itica, su desarrollo de Taylor alrededor de $z=a$ es
\begin{equation}
  \label{eq:taylor_en_a}
  f(z)=f(a)+\eval{\dv{}{z}f(z)}_{z=a}(z-a)+\dots,
\end{equation}
de modo que en \eqref{eq:F_cociente} el factor $(z-a)$ se cancela y
\begin{equation}
  \label{eq:F_analitica}
  F(z)=a_1+a_2(z-a)+\dots
\end{equation}
es anal\'itica en todo $D$, incluido $z=a$. Por el teorema de Cauchy
\eqref{eq:integral_cerrada_nula} su integral sobre un contorno cerrado se
anula:
\begin{equation}
  \label{eq:derivacion_fic}
  \begin{aligned}
    0&=\oint_{\ga}\frac{f(z)-f(a)}{z-a}\dd z\\
     &=\oint_{\ga}\frac{f(z)}{z-a}\dd z
       -f(a)\oint_{\ga}\frac{\dd z}{z-a}\\
     &=\oint_{\ga}\frac{f(z)}{z-a}\dd z-2\pi i f(a),
  \end{aligned}
\end{equation}
usando \eqref{eq:integral_fundamental} en el \'ultimo paso.

\begin{teorema}[F\'ormula integral de Cauchy]
  \label{teo:fic}
  \begin{equation}
    \label{eq:fic}
    f(a)=\frac{1}{2\pi i}\oint_{\ga}\frac{f(z)}{z-a}\dd z,
  \end{equation}
  o, escribiendo la variable de integraci\'on como $\xii$,
  \begin{equation}
    \label{eq:fic_xi}
    f(z)=\frac{1}{2\pi i}\oint_{\ga}\frac{f(\xii)}{\xii-z}\dd\xii .
  \end{equation}
\end{teorema}

Es decir, los valores de una funci\'on anal\'itica en el interior de $\ga$
quedan determinados por sus valores sobre $\ga$.

\begin{ejemplo}
  \label{ej:fic_aplicada}
  Si $g(z)=e^{2z^2}-\cos^3\qty{z^{1/2}+3z^{1/3}}$ es anal\'itica dentro de
  $\ga$ y $z=2$ est\'a en el interior, entonces
  \begin{equation}
    \label{eq:fic_aplicada}
    \oint_{\ga}\frac{g(z)}{z-2}\dd z=2\pi i\,g(2),
  \end{equation}
  sin necesidad de evaluar ninguna integral.
\end{ejemplo}

\begin{definicion}[\'Indice]
  \label{def:indice}
  \begin{equation}
    \label{eq:indice}
    I=\frac{1}{2\pi i}\oint_{\ga}\frac{\dd z}{z-a}\in\Z,
  \end{equation}
  el n\'umero de vueltas que $\ga$ da alrededor de $a$.
\end{definicion}

\subsection{F\'ormula de Cauchy generalizada}
\label{subsec:cauchy_generalizada}

Derivando \eqref{eq:fic_xi} bajo el signo integral,
\begin{equation}
  \label{eq:fic_derivada}
  \dv{}{z}f(z)=\dv{}{z}\set{\frac{1}{2\pi i}
    \oint_{\ga}\frac{f(\xii)}{\xii-z}\dd\xii}
    =\frac{1}{2\pi i}\oint_{\ga}\frac{f(\xii)}{(\xii-z)^2}\dd\xii,
\end{equation}
y en general
\begin{equation}
  \label{eq:fic_general}
  \dnv{}{z}{n}f(z)
    =\frac{n!}{2\pi i}\oint_{\ga}
      \frac{f(\xii)}{(\xii-z)^{n+1}}\dd\xii .
\end{equation}

\subsection{Desigualdad de Cauchy y teorema de Liouville}
\label{subsec:liouville}

Combinando Taylor \eqref{eq:taylor_complejo} con \eqref{eq:fic_general}, los
coeficientes del desarrollo $f(z)=\sum_n C_n(z-a)^n$ son
\begin{equation}
  \label{eq:coeficientes_fic}
  C_n=\frac{1}{2\pi i}\oint_{\ga}
    \frac{f(\xii)}{(\xii-a)^{n+1}}\dd\xii .
\end{equation}
Acotando con \eqref{eq:cota_integral} y deformando $\ga$ a la circunferencia
$\xii=a+\rho e^{i\vph}$, para la cual $\abs{\dd\xii}=\rho\dd\vph$,
\begin{equation}
  \label{eq:cota_coeficientes}
  \begin{aligned}
    \abs{C_n}
      &\le\frac{1}{2\pi}\oint_{\ga}
        \frac{\abs{f(\xii)}\abs{\dd\xii}}{\abs{\xii-a}^{n+1}}
       =\frac{1}{2\pi}\int_{0}^{2\pi}
        \frac{\abs{f(\vph)}\rho\dd\vph}{\rho^{n+1}}\\
      &=\frac{1}{2\pi}\int_{0}^{2\pi}
        \frac{\abs{f(\vph)}}{\rho^{n}}\dd\vph .
  \end{aligned}
\end{equation}
Si $f$ est\'a acotada, $\abs{f}\le M$, se obtiene la
\textbf{desigualdad de Cauchy}
\begin{equation}
  \label{eq:desigualdad_cauchy}
  \abs{C_n}\le\frac{M}{\rho^{n}} .
\end{equation}
Como $0\le\rho<R$ con $R$ el radio de convergencia, $1/R\le 1/\rho$ y el
supremo se alcanza en
\begin{equation}
  \label{eq:supremo_coeficientes}
  \abs{C_n}\le\frac{M}{R^{n}} .
\end{equation}

\begin{teorema}[Liouville]
  \label{teo:liouville}
  Toda funci\'on anal\'itica y acotada en todo el plano debe ser constante.
\end{teorema}

\begin{demostracion}
  Si $f$ es entera y acotada, $R$ puede tomarse arbitrariamente grande en
  \eqref{eq:supremo_coeficientes}, de modo que $C_n=0$ para todo $n\ge 1$ y
  s\'olo sobrevive $C_0$.
\end{demostracion}

\subsection{Singularidades removibles}
\label{subsec:singularidades_removibles}

\begin{teorema}
  \label{teo:singularidad_removible}
  Supongamos $f(z)$ anal\'itica en $D$ excepto en un punto aislado $z=a$. Una
  condici\'on necesaria y suficiente para que exista una extensi\'on
  anal\'itica de $f$ a todo $D$, incluido $z=a$, es
  \begin{equation}
    \label{eq:singularidad_removible}
    \lim_{z\rightarrow a}(z-a)f(z)=0 .
  \end{equation}
  Adem\'as, la funci\'on extendida es \'unica.
\end{teorema}

\begin{ejemplo}
  \label{ej:sinc}
  Sea $f(z)=\dfrac{\sin z}{z}$. Como
  $\lim_{z\rightarrow 0}\dfrac{\sin z}{z}=1$, se tiene
  \begin{equation}
    \label{eq:sinc}
    \lim_{z\rightarrow 0}zf(z)
      =\lim_{z\rightarrow 0}\sin z=0,
  \end{equation}
  y la singularidad en $z=0$ es removible.
\end{ejemplo}

\subsubsection{Construcci\'on iterativa de Taylor}
\label{subsubsec:construccion_taylor}

La idea anterior se puede iterar. Partiendo de
\begin{equation}
  \label{eq:F0}
  F_0(z)=F(z)=\frac{f(z)-f(a)}{z-a},
\end{equation}
que cumple $\lim_{z\rightarrow a}F(z)=\eval{\dv{}{z}f(z)}_{z=a}$ y por lo
tanto $\lim_{z\rightarrow a}(z-a)F(z)=0$, definimos recursivamente
\begin{equation}
  \label{eq:Fn_recursion}
  F_n(z)=\frac{F_{n-1}(z)-F_{n-1}(a)}{z-a},
\end{equation}
de modo que
\begin{equation}
  \label{eq:limite_Fn}
  \lim_{z\rightarrow a}F_1(z)
    =\eval{\dv{}{z}F(z)}_{z=a}
    =\frac{1}{2}\eval{\ddv{}{z}f(z)}_{z=a},
\end{equation}
y en general
\begin{equation}
  \label{eq:Fn_derivada}
  F_n(a)=\frac{1}{n!}\eval{\dnv{}{z}{n}f(z)}_{z=a} .
\end{equation}
Despejando hacia atr\'as la recursi\'on \eqref{eq:Fn_recursion} se reconstruye
el desarrollo de Taylor con residuo.

\begin{teorema}
  \label{teo:taylor_residuo}
  Si $f(z)$ es anal\'itica en una regi\'on $D$ que contiene a $z=a$, entonces
  \begin{equation}
    \label{eq:taylor_residuo}
    \small
    \begin{aligned}
      f(z)&=f(a)+\frac{1}{1!}f'(a)(z-a)
        +\frac{1}{2!}f''(a)(z-a)^2+\dots\\
        &\quad+\frac{f^{(n)}(a)}{n!}(z-a)^n
        +F_{n+1}(z)(z-a)^{n+1},
    \end{aligned}
  \end{equation}
  con $F_{n+1}(z)$ anal\'itica en $D$.
\end{teorema}

\subsection{Ceros}
\label{subsec:ceros}

Supongamos que $f(z)$ es anal\'itica en $D$ y que $f(a)=0$. Si adem\'as
\begin{equation}
  \label{eq:cero_orden_n}
  \eval{\dnv{}{z}{\ka}f(z)}_{z=a}=0
  \quad\forall\,0\le\ka\le n-1,
  \qquad
  \eval{\dnv{}{z}{n}f(z)}_{z=a}\neq 0,
\end{equation}
decimos que $z=a$ es un \textbf{cero de orden $n$}. En tal caso el desarrollo
de Taylor empieza en la potencia $n$ y podemos factorizar
\begin{equation}
  \label{eq:factorizacion_cero}
  f(z)=(z-a)^n\tilde{f}(z),
  \qquad \tilde{f}(a)\neq 0,
\end{equation}
con $\tilde{f}$ anal\'itica. Si hay varios ceros $\set{z_j}$ en $D$,
\begin{equation}
  \label{eq:factorizacion_multiple}
  f(z)=(z-z_1)^{m_1}(z-z_2)^{m_2}\dots(z-z_r)^{m_r}\tilde{f}(z),
\end{equation}
donde $m_k$ es la \textbf{multiplicidad} de cada cero y
$\sum_{k=1}^{r}m_k=m$ es el orden total.

\subsection{Polos y series de Laurent}
\label{subsec:polos_laurent}

Sea $f(z)$ anal\'itica en $D$ excepto en un conjunto de puntos aislados. En
cada uno de ellos se distinguen dos casos:
\begin{itemize}
  \item si $\lim_{z\rightarrow a}(z-a)f(z)=0$, la singularidad es
        \textbf{removible};
  \item si $\lim_{z\rightarrow a}f(z)=\infty$, tenemos un \textbf{polo}.
\end{itemize}

Supongamos que $a$ es un polo y tomemos un disco de radio $\tilde{R}<R$ donde
$f(z)\neq 0$ para $0<\abs{z-a}<\tilde{R}$. Definimos
\begin{equation}
  \label{eq:g_reciproca}
  g(z)=\frac{1}{f(z)},
\end{equation}
anal\'itica salvo en $z=a$ y con $\lim_{z\rightarrow a}g(z)=0$; por el
Teorema~\ref{teo:singularidad_removible} se extiende anal\'iticamente al
disco. Entonces $z=a$ es un cero de orden $n$ de $g$, y por
\eqref{eq:factorizacion_cero}
\begin{equation}
  \label{eq:g_factorizada}
  g(z)=(z-a)^n\tilde{g}(z),
  \qquad \tilde{g}(a)\neq 0 .
\end{equation}
En consecuencia
\begin{equation}
  \label{eq:f_polo}
  f(z)=\frac{1}{(z-a)^n}\tilde{f}(z),
  \qquad
  \tilde{f}(z)=\frac{1}{\tilde{g}(z)} .
\end{equation}
Como $\tilde{f}$ es anal\'itica, admite desarrollo de Taylor, y sustituyendo
\begin{equation}
  \label{eq:laurent_desarrollo}
  \begin{aligned}
    f(z)&=\frac{1}{(z-a)^n}\sum_{k=0}^{\infty}C_k(z-a)^k\\
        &=\frac{B_n}{(z-a)^n}+\dots+\frac{B_1}{(z-a)}
          +B_0+B_1'(z-a)+\dots
  \end{aligned}
\end{equation}

\begin{definicion}[Serie de Laurent]
  \label{def:laurent}
  En general, alrededor de una singularidad de $f(z)$ escribimos
  \begin{equation}
    \label{eq:laurent}
    f(z)=\sum_{k=-\infty}^{\infty}C_k(z-a)^k,
  \end{equation}
  que se separa en
  \begin{equation}
    \label{eq:laurent_partes}
    f(z)=\underbrace{\sum_{k=-\infty}^{-1}C_k(z-a)^k}
           _{\text{parte principal o singular}}
      +\underbrace{\sum_{k=0}^{\infty}C_k(z-a)^k}
           _{\text{parte anal\'itica}} .
  \end{equation}
\end{definicion}

Seg\'un cu\'antos t\'erminos tenga la parte principal:
\begin{itemize}
  \item \textbf{polo de orden $n$} si tiene exactamente $n$ t\'erminos;
  \item \textbf{singularidad esencial} si tiene infinitos t\'erminos.
\end{itemize}


\end{document}
